
\label{ch:discussion}

\begin{ataglance}[How to read this chapter]
The empirical chapters return numbers. This chapter says what the numbers would
mean, and what they would still leave open.

It runs in seven moves and a closing section. The first states the principal
findings against the four aims and marks which of them depend on the analysis and
which do not. The second fixes, before any result is seen, how each possible
pattern of results will be read, so that the reading cannot drift towards
whatever the numbers turn out to be. The third sets the $H_{\text{potential}}$
against the measures it would have to sit beside: capability theory, frailty
indices, health utilities, the social determinants literature, and the standards
now used to judge prediction models. The fourth separates what the design earns
from what it cannot earn, since several strengths and their matching limitations
come from one decision, and it carries forward the shortcomings that the checks of
the data against their source documents left unresolved. The fifth asks what
follows for measurement, for method, and for clinical and policy use, holding each
claim to what a single development cohort can support. The sixth reads the thesis
as a systems problem: what kind of whole health capacity is, where a change does
most, and what happens to a score once people act on it. The seventh sets out what
would have to happen next, external validation first. Two closing sections follow.
The first states what the thesis contributes, separating what stands now from what
waits on the NHATS results. The second states what the chapter settles and what it
passes to Chapter~\ref{ch:conclusions}.

Two conventions run throughout. Each statistical term is given with its clinical
or laboratory equivalent at first use, and the chapter says where the parallel
stops holding. Claims about performance are restated for the reference
profiles of Table~\ref{tab:profiles} and are written conditionally, because the
numbers are not yet in.
\end{ataglance}

\begin{quote}
\small\itshape
Note to the reader. Sections that interpret specific results are marked with placeholders. These are completed after the analysis of Chapters~\ref{ch:results1} and~\ref{ch:results2}. The strengths, limitations, and future-work sections stand on their own. This note is removed in the final version.
\end{quote}

\noindent This chapter answers one question. When the analyses of
Chapters~\ref{ch:results1} and~\ref{ch:results2} are complete, what will they
entitle anyone to say about health, about how it is measured and about what a
service should do? The question has to be answered now, before the numbers exist,
for a plain reason. A discussion written after the results is under pressure to
find that whatever happened was what the theory expected. A strong loading becomes
support for the framework; a weak one becomes a fault of the survey. Either reading
may be right in a given case, but if both are open once the numbers are known,
neither tests anything. Writing the interpretation first, and giving each possible
outcome its reading in advance, removes that freedom. The chapter is therefore
conditional wherever it concerns findings, and firm only where it concerns the
design, the data and the argument, which already exist.

\noindent The thesis set out to develop, validate and evaluate a measure of health, the $H_{\text{potential}}$, that treats health as a person's latent capacity to maintain health over time, not as current function or accumulated disease. Existing measures each capture part of health. Measures of disease burden describe illness, quality-of-life instruments describe perceived health \citep{whitehead_2010}, frailty indices count physiological deficits, and quality-adjusted life years supply utilities for allocating resources. None was designed to bring biological, functional, social and structural determinants into one construct that represents future capability.

The gap matters more as populations age and chronic disease comes to dominate morbidity. Health systems now care for people who live for decades with several chronic conditions. Understanding their health requires knowing their capacity to keep function, adapt, use the care available and recover from biological, psychological and social shocks. Frameworks of healthy ageing, the capability approach and the social determinants of health all stress these dimensions, but they have rarely been brought into one measurement framework that can be validated and tested against outcomes.

The $H_{\text{potential}}$ was proposed to fill that gap. It treats health as a latent capability arising from three domains. Biomedical Potential is the physiological and functional reserve that lets a person stay well despite ageing and disease. Social Modifiers are the interpersonal, behavioural and community resources that decide whether that reserve turns into good outcomes. Systemic Burdens are the healthcare, organisational and structural conditions that help or hinder the realisation of health. Together the domains reach beyond a biomedical model while keeping enough structure to be measured.

The thesis ran in sequence from theory to measurement, from measurement to prediction, and from prediction to translation, on the principle that a construct should be defined before it is put into statistical form and tested. Conceptual development, psychometric validation, prognostic assessment and disease-specific adaptation were kept separate throughout, which reduces the risk of circular reasoning and makes the result easier to interpret.

Chapter~\ref{ch:theory} established the theoretical foundations of the $H_{\text{potential}}$. Drawing on the capability approach of \citet{sen_1993}, the World Health Organization's concept of healthy ageing, contemporary frailty theory and the literature on social determinants of health, the chapter argued that health should be understood as a dynamic capability rather than solely as a current health state. Existing measures were critically examined, highlighting their considerable strengths while identifying a common limitation: each captures an important dimension of health, but none explicitly integrates the biological, social and systemic determinants that jointly influence future health trajectories. Chapter~\ref{ch:methods} translated this conceptual framework into a reproducible empirical methodology, developing a two-stage analytical strategy in which latent variable estimation and prognostic evaluation were performed as separate stages rather than simultaneously. The empirical analyses of Chapters~\ref{ch:results1} and~\ref{ch:results2} evaluated the psychometric and prognostic properties of the framework within the National Health and Aging Trends Study (NHATS), and the final empirical chapter, Chapter~\ref{ch:translational}, extended the work by proposing a structured pathway for translation to haemophilia. Collectively, these chapters address the four aims of the study, as summarised in Figure~\ref{fig:contributions}.

\begin{figure}[htbp]
\centering
\begin{tikzpicture}[node distance=5mm]
\node[hpbox, text width=2.75cm] (a1) {\textbf{Aim 1}\\Theory\\[1pt]\footnotesize Chapter~\ref{ch:theory}};
\node[hpbox, text width=2.75cm, right=of a1] (a2) {\textbf{Aim 2}\\Measurement model\\[1pt]\footnotesize Chapter~\ref{ch:results1}};
\node[hpbox, text width=2.75cm, right=9mm of a2] (a3) {\textbf{Aim 3}\\Prognostic evaluation\\[1pt]\footnotesize Chapter~\ref{ch:results2}};
\node[hpbox, text width=2.75cm, right=of a3] (a4) {\textbf{Aim 4}\\Translational mapping\\[1pt]\footnotesize Chapter~\ref{ch:translational}};
\draw[hparrow] (a1) -- (a2);
\draw[hparrow] (a2) -- (a3);
\draw[hparrow] (a3) -- (a4);
\coordinate (mid) at ($(a2.east)!0.5!(a3.west)$);
\draw[hpdash] ($(mid)+(0,1.25)$) -- ($(mid)+(0,-1.25)$);
\node[hpnote, text width=6.2cm, below=1.35cm of mid] {Measurement is fixed here, before any outcome is used};
\end{tikzpicture}

\captionsetup{justification=justified, width=\linewidth}
\caption{The sequential structure of the thesis across the four study aims.}
\label{fig:contributions}
\end{figure}

% =====================================================================
\section{Summary of Principal Findings}
\label{sec:disc-findings}
% =====================================================================

\noindent A thesis with four aims can fail in four ways, and the failures do not
cancel. A construct can be coherent and still useless, or useful in one cohort
and empty in another. Reporting the findings aim by aim keeps those verdicts
apart, so that a weak result in one stage is not paid for with a strong result
in the next.

\noindent Each stage of the programme answered its own question and built on the one before, so the contribution of the thesis rests on the four stages together, not on a single statistical result.

The first principal finding concerns the theoretical development of the $H_{\text{potential}}$. The thesis demonstrated that a coherent conceptual framework can be constructed by integrating ideas drawn from the capability approach, healthy ageing, frailty science and the social determinants of health. Rather than treating biological disease, social environment and healthcare systems as independent explanatory variables, the proposed framework positions them as interacting determinants of a common latent capability. Conceptually, this represents a shift away from describing current health status towards measuring an individual's future capacity to maintain health.

The second principal finding relates to measurement of the latent construct.
Two questions decide whether it stands. The first is whether the indicators hold
together as three domains at all. The second is discriminant validity, which
asks whether two instruments are measuring the same thing under different names.
A score that correlates almost perfectly with a frailty index has passed no test
worth passing, because it adds nothing the frailty index does not already
supply.

A laboratory reader already knows this sequence from assay validation. A latent
domain cannot be observed directly. It is estimated from several survey
indicators, much as an analyte concentration is inferred from an assay signal
through a calibration curve. The measurement model is that calibration. Its
loadings say how strongly each indicator tracks the domain, as one would ask how
strongly a given assay tracks the analyte it is meant to measure. Reliability is
the counterpart of precision: whether repeated or parallel readings agree.
Validity is the counterpart of accuracy: whether the readings agree with the
quantity they claim to measure. And discriminant validity is the question a
laboratory asks of a new assay that correlates closely with an existing one: if
the two always agree, the new one adds work and no information. The parallel has
one limit worth stating at once. An assay can be checked against a reference
material of known concentration. No reference standard of health capacity exists,
so the measurement model can be checked only for internal coherence and for its
relation to other measures and to later events.

\medskip
\noindent\emph{[Placeholder: summarise the psychometric findings from Chapter~\ref{ch:results1}. Report whether the hypothesised three-domain structure demonstrated satisfactory model fit, whether reliability and construct validity supported the proposed measurement model, whether measurement invariance \citep{meredith_1993,vandenberg_lance_2000} was established where tested, and whether the resulting factor scores provided a stable representation of the $H_{\text{potential}}$. Interpretation should focus on what these findings imply for the feasibility of measuring latent health capability.]}

The third principal finding concerns prognostic performance. More than one
quantity carries that judgement, and each has a plain clinical form.
Discrimination is the ability to rank patients by risk, as a diagnostic test
ranks by likelihood: given one person who has the event and one who outlives
them, how often the score puts the two in the right order. Calibration asks whether a predicted
10 per cent risk means that about one in ten such people had the event, which is
what matters the moment a number is spoken aloud in a consultation. Incremental
value asks whether the new test changes management over the one already in use.

Each has a laboratory counterpart. Discrimination for a survival outcome is
summarised by the $C$-statistic, which reads like the area under an ROC curve: 0.5
is no better than chance and 1 is perfect ordering. Calibration corresponds to
checking an assay against standards of known value, where an assay that ranks
samples correctly but reads every one too high still fails. Incremental
value corresponds to the test a laboratory applies before adding a new assay to a
panel. A new assay that correlates with an existing one adds value only if it
changes decisions that the existing one would have got wrong. A hazard ratio,
familiar from trials, is the ratio of event rates between two groups at each
moment of follow-up; here it compares people one standard deviation apart on the
score.

\medskip
\noindent\emph{[Placeholder: summarise the findings from Chapter~\ref{ch:results2}. Report whether the $H_{\text{potential}}$ independently predicted each co-primary outcome, five-year all-cause mortality and, separately, first overnight hospital stay, the magnitude of each association, the discrimination and calibration achieved, and whether the framework demonstrated incremental value over existing approaches such as frailty indices, QALY proxies \citep{whitehead_2010,drummond2015} and comorbidity measures. Interpretation should emphasise clinical significance rather than statistical significance alone.]}

\begin{reading}[A hazard ratio on the absolute scale]
Suppose the analysis returns a hazard ratio per standard deviation of the
$H_{\text{potential}}$. On its own that number says little about absolute risk,
because it is a ratio of rates. The conversion needs a baseline. If five-year
mortality at the mean of the score is near one in five, then a hazard ratio of,
say, $1.4$ applied to the standard deviation by which Profile~A sits below the
mean moves its risk to about $27$ per cent: roughly seven extra deaths per
hundred people at that profile over five years, or one extra death for every
fifteen.

Whether that is worth acting on depends on what action exists. If most of the
shortfall sits in the systemic domain, the hazard ratio is in part a statement
about structural conditions. If it sits in the biomedical domain, it is a
statement about physiology. The same ratio carries different instructions depending on
which domain produced it, and the multiplicative form is what allows the score to
be taken apart that way.
\end{reading}

The fourth principal finding concerns translation beyond the developmental population. The thesis demonstrated that adaptation of the $H_{\text{potential}}$ to haemophilia can be undertaken without altering the underlying conceptual framework. Instead of constructing an entirely new disease-specific measure, the proposed translational strategy preserves the latent domains while replacing their observable indicators with clinically relevant haemophilia-specific variables. The chapter further established a structured methodology comprising expert consensus, psychometric re-estimation, recalibration and external validation before any clinical implementation. Although empirical validation within a haemophilia cohort lies beyond the scope of the present study, the translational framework provides a clear and scientifically defensible pathway for future investigation.

Taken together, these findings suggest that the contribution of the present thesis extends beyond development of another composite health index. The work proposes an alternative way of conceptualising health itself. Instead of measuring health as a static collection of diseases, symptoms or functional limitations, the $H_{\text{potential}}$ conceptualises health as a dynamic, multidimensional capability that emerges through interactions among biological reserve, social resources and healthcare systems. Importantly, the overall conclusions of the thesis remain appropriately bounded. The psychometric and prognostic findings are derived from a single development cohort and therefore require confirmation through independent external validation before routine application can be recommended. Likewise, the haemophilia framework presented in Chapter~\ref{ch:translational} represents a conceptual translation rather than a completed disease-specific measurement model.

Of the four findings, the first and the fourth do not wait on data. The
theoretical framework and the haemophilia mapping are arguments, and they stand or
fall on their reasoning. The second and the third do wait, and they are not
independent of each other. A prognostic result for a construct whose measurement
model failed would be a result for an unknown quantity: a weighted sum of eleven
indicators that predicts something, but not evidence about capability. The order in
which the two blanks are filled therefore matters, and so do the rules for filling
them.

% =====================================================================
\section{How the Results Will Be Read}
\label{sec:disc-reading}
% =====================================================================

\noindent The previous section left two findings blank and showed that the second
depends on the first. What it did not settle is how either blank will be read once
it is filled. This section settles that now. Each rule below names a pattern the
analysis could return and the reading that pattern will receive, so that the
reading is fixed before the result is seen.

\subsection{The Order of Reading}
\label{sec:disc-read-order}

The measurement results are read first and decide what the prognostic results are
about. If the three-domain model fits, its loadings are adequate and its domains
are distinct from the frailty comparator, the score passed to the survival models
is an estimate of the construct the theory describes, and a prognostic gain is
evidence about that construct. If the model fails, the survival models still run
and may still return a useful predictor, but the predictor is then a composite
whose meaning has not been shown. Reading the prognostic results without first
asking which of these two cases holds would let a good hazard ratio vouch for a
measurement model that did not earn it. A laboratory follows the same order for
the same reason. A patient result from a run whose calibration or quality control
failed is not reported, however plausible it looks, because nothing then ties the
number to the analyte. Here the prognostic results are still reported in that
case, but as results about an unnamed composite, not about health capacity.

Within the prognostic results the order follows from use. A score used only to
rank people, to decide who is visited first, needs discrimination and nothing
else, since a ranking does not depend on the scale of the predicted risks. A score
used to act at a threshold, to offer an assessment to everyone above 20 per cent,
needs calibration as well, because a model that ranks well but overstates every
risk sends too many people across the threshold. Net benefit is read last, since
it depends on both and on the threshold chosen \citep{vickers_elkin_2006}.
Incremental value is read against the comparator in use, in the paired form of
Derivation~\ref{der:incremental}, and in absolute risks as in Worked
Example~\ref{ex:disc:incremental}, because a ratio alone cannot show whether a
decision changes. The last question is whether the non-compensatory realisation
model of Chapter~\ref{chap:crf} does better than a plain sum of the same domains.
It comes last because it asks how the domains combine, which presumes that they
have been measured and that the score predicts. The four steps, and the branch
taken under each pattern, are drawn together in Figure~\ref{fig:disc-decision} at
the end of this section.

\subsection{Reading the Measurement Results}
\label{sec:disc-read-measure}

The measurement model can return one of four broad patterns, and each has its own
reading.

In the first, the three-domain structure fits, every domain is defined by
indicators that load well, and the correlations between domains are high enough to
support a common higher-order construct but low enough that the domains remain
distinct. That is the pattern the theory predicts, and it would show that latent
health capability can be measured in this cohort with these indicators. It would
not show more than that.

In the second, the biomedical domain is well defined and the social and systemic
domains are weak. Two readings are open, and the analysis plan contains the test
that separates them. The weakness may lie in the indicators: NHATS was not designed
to measure social enablement or structural freedom, and a domain measured by two or
three distant proxies will load poorly whatever its importance. Or the weakness may
lie in the specification: social and structural conditions may not be reflections
of one underlying quantity but components that combine to form it, in which case a
reflective factor model is the wrong tool. The distinction has a laboratory form.
In a reflective model the indicators are several readings of one underlying
quantity, as several assays of the same analyte should rise and fall together. In
a formative model the indicators are separate components that together make up
the quantity, as a panel score is assembled from tests that measure different
things and need not agree with one another. Asking components of the second kind
to behave like readings of the first kind will make them look poor even when each
is accurate. The formative and hybrid specifications
among the sensitivity analyses of Chapter~\ref{ch:methods} test the second reading.
If those specifications behave better, the conclusion is that the contextual
domains are composites, and the framework survives with a different measurement
model. If they behave no better, the conclusion is that NHATS does not measure these
domains well enough, and the framework is untested on them rather than refuted.

In the third, the three domains cannot be told apart: their correlations approach
one and a single factor fits as well as three. The score may still predict
outcomes. But the realisation function of Chapter~\ref{chap:crf} treats social
enablement and structural freedom as separate inputs whose shortfalls combine, and
if the data cannot separate them, that function has no empirical footing in this
cohort. The thesis would then report a single-factor health score and set the
three-domain theory aside as unsupported by these data.

In the fourth, the structure fits in the cohort as a whole but measurement
invariance \citep{meredith_1993,vandenberg_lance_2000} fails across sex or across
age band. The score could then be used within a group but not compared across
groups, and every subgroup comparison in Chapter~\ref{ch:results2} would carry that
qualification. Measurement invariance asks whether the instrument gives the same
reading for the same level of capacity in men and women, or in younger and older
respondents, as an assay should give the same reading for the same concentration
in different patient groups or different laboratories. Where it does not, results
can be compared within a group, as results can be compared within one laboratory,
but a difference between groups may belong to the instrument and not to the
people.

Two of the stability checks reported in Chapter~\ref{ch:results1} bear on all four
patterns. They play the part that a change of method or reagent lot plays in assay
validation: a result that holds only under one estimation method is a property of
the method. If the loadings move materially between the maximum-likelihood and WLSMV
solutions (two standard ways of fitting the model, the second built for answers on
ordered categories), or when single indicators are removed in turn, the structure depends on
a modelling choice or on one item, and is reported as fragile. If the loadings and
biomedical scores move materially when the cognition of people whose proxies said
they could not be tested is set to the lowest possible score in place of the imputed
value, the biomedical domain depends on an assumption about those people that the
data cannot check (Appendix~\ref{app:datachecks}). That result is then reported as
a range between the two rules, not as one value.

\subsection{Reading the Prognostic Results}
\label{sec:disc-read-prog}

The prognostic results raise three questions that a single performance figure
would hide.

The first is whether the two co-primary outcomes agree. Five-year mortality and the
first overnight hospital stay are reported separately, not combined, and they may
part company. A hospital stay is partly an event of illness and partly an event of
access: to be admitted, a person has to reach a hospital that will admit them. A
person with low structural freedom may therefore be at higher risk of death and, at
the same time, at no higher or even lower recorded risk of admission. If the score
predicts death well and admission poorly, that pattern is one the framework allows
for, and it is read as a statement about what an admission measures, not as a
failure of the score. If the score predicts admission well and death poorly, the
reading is less comfortable, because it suggests the score tracks contact with
services more than the capacity to stay alive. Death also competes with admission:
a person who dies cannot then be admitted, so a score that predicts death strongly
can appear to predict admission less well for that reason alone. The competing-risk
treatment of Chapter~\ref{ch:methods} is there to keep these two effects apart. A
competing risk is an event that prevents the event of interest from being seen,
as death does here; it differs from censoring, the cohort-study case of a patient
lost to follow-up, who might still have had the event unobserved.

The second is whether any gain over the frailty comparator survives the removal of
indicators that overlap with it. One sensitivity analysis drops those indicators.
If the gain persists, it comes from the domains the frailty index does not cover.
If it vanishes, the $H_{\text{potential}}$ was adding a second copy of the frailty
index's own information, and the gain is reported as duplication.

The third is whether the results depend on the people placed at the floor of the
domain scale. The preregistered rule maps each domain onto $[0,1]$ using the 1st
and 99th percentiles of the development split and clips values outside them, so
that anyone below the 1st percentile on a domain is placed at, in effect, zero.
Under the geometric realisation function a zero on one domain sets the whole score
to zero whatever the other two are. The clipping resembles reporting every result
below the lower limit of quantitation as that limit: the values become alike
whatever their true level. The parallel stops there, because a low result for one
analyte does not erase the results for the others, while under the geometric form
a zero on one domain does. The fixed-floor analysis maps the same bounds
onto $[0.05,1]$ instead. If the conclusions are the same under both rules, the
floor does not drive them. If they differ, a small group at the bottom of one
domain is carrying the result, and the report says so.

Table~\ref{tab:disc-patterns} gathers these readings in one place. Each row names
a pattern the analysis could return, the reading it will receive, and what it
would mean for the next step of the programme. The table records commitments, not
results: no row has yet been observed.

\begin{table}[htbp]
\centering
\caption{Patterns the analysis could return and the reading fixed for each in advance. No row is a result.}
\label{tab:disc-patterns}
\footnotesize
\renewcommand{\arraystretch}{1.15}
\begin{tabular}{p{4.6cm} p{5.2cm} p{4.0cm}}
\toprule
\textbf{Pattern} & \textbf{Reading} & \textbf{Consequence} \\
\midrule
Three domains fit, distinct from frailty & Construct measurable in this cohort & Proceed to external validation \\
Biomedical strong, contextual domains weak; formative form better & Contextual domains are composites & Revise measurement model, keep theory \\
Biomedical strong, contextual domains weak; formative form no better & NHATS measures these domains poorly & Seek a cohort with richer indicators \\
Domains inseparable & Single-factor score only & Set aside three-domain theory for these data \\
Invariance fails across sex or age & Within-group use only & Qualify all subgroup comparisons \\
Gain over frailty lost when overlap removed & Duplication, not new information & Report as null for incremental value; counts against the claim of new information \\
Death predicted, admission not & Admission partly measures access & Report outcomes separately \\
Conclusions change with the $[0.05,1]$ floor & Result carried by lowest scorers & Report both; review the scaling rule \\
\bottomrule
\end{tabular}
\end{table}

\subsection{What Would Count Against the Framework}
\label{sec:disc-read-against}

A framework that every result could be read as supporting would make no claim. It
is therefore worth stating which results would count against this one. Three would.
The first is a measurement model in which the domains cannot be separated, since
the theory's main claim is that health capacity has three distinct sources. The
second is a score that adds nothing to the frailty index at any decision threshold
once overlapping indicators are removed, since the case for a broader construct
rests on the information the broader domains add. The third is a score that remains
badly calibrated after internal validation, since a capacity measure that misstates
absolute risk cannot be used at a threshold. A fourth test sits in
Chapter~\ref{chap:crf}, whose pre-specified head-to-head protocol
(Section~\ref{sec:E-falsification}) sets the non-compensatory realisation function
against an additive composite of the same domains and against the factor score. If
the realisation function cannot be told apart from the additive composite, or does
worse on calibration and prediction error, that chapter's own rule applies: the
empirical claim for the realisation function is narrowed to a theoretical model, and
the factor score is kept as the measurement backbone.

Figure~\ref{fig:disc-decision} sets these commitments out in the order they will
be read. It is a reading plan, not a result, and holds nothing that the text and
Table~\ref{tab:disc-patterns} do not already state. Each step on the left is read
only after the one above it. Each row records a pattern the analysis could
return and the reading this chapter has already fixed for it. The line style and
shading of a branch give its consequence for the framework: a solid line where the
framework is kept, a dashed line where it is kept with a stated qualification, and
a dotted line on a darker ground where the result counts against the framework or
narrows one of its claims. Three features of the figure carry the argument. The
measurement step comes first and has the most branches, because every later step
is about the construct only if this one holds. Two of the patterns that count
against the framework sit in the comparison step, where a score can pass every
test of its own and still add nothing to the frailty index already in use. And
the realisation model has its own failure branch, which narrows the claim for that
model while leaving the factor score in place, so a result against the
realisation model is not a result against the measure. Whatever branch the
analysis takes, the next step is the same: external validation in a second
cohort.

\begin{figure}[htbp]
\centering
% Decision diagram for reading the results, Chapter 9 (Discussion).
% Greyscale and line style carry every distinction: solid = framework kept,
% dashed = kept with a qualification, dotted = counts against or narrows.
{\hpdiagramsizes
\begin{tikzpicture}[
  pat/.style={draw, rounded corners=2pt, align=flush left, font=\footnotesize,
              text width=4.3cm, inner sep=3pt, minimum height=0.74cm},
  con/.style={pat, text width=4.9cm},
  kept/.style={fill=black!6, thick, solid},
  qual/.style={fill=white, thick, dashed},
  agst/.style={fill=black!20, very thick, dotted},
  stepbox/.style={hpkey, text width=2.3cm, font=\footnotesize, align=center},
  link/.style={-{Stealth[length=1.6mm]}, semithick},
  hd/.style={font=\footnotesize\bfseries, align=left},
]
\def\xs{1.35}   % step column centre
\def\xp{5.15}   % pattern column centre
\def\xc{10.75}  % conclusion column centre
\def\xe{2.65}   % step column right edge
% Column headings
\node[hd] at (\xs,0.75) {Step};
\node[hd, anchor=west] at (\xp-2.3,0.75) {Pattern the analysis returns};
\node[hd, anchor=west] at (\xc-2.6,0.75) {Reading fixed in advance};
% One branch: #1 y, #2 style, #3 pattern, #4 conclusion
\def\hpdbranch#1#2#3#4{%
  \node[pat,#2] (p) at (\xp,#1) {#3};
  \node[con,#2] (c) at (\xc,#1) {#4};
  \draw[link] (\xe,#1) -- (p.west);
  \draw[link,#2,fill=none] (p.east) -- (c.west);
}
% ---- Step 1: measurement
\hpdbranch{0}{kept}{Three domains fit, load well and stay distinct from frailty}{Construct measurable in this cohort; step 2 is then evidence about it}
\hpdbranch{-0.92}{qual}{Contextual domains weak; formative form does better}{Composites: revise the measurement model, keep the theory}
\hpdbranch{-1.84}{qual}{Contextual domains weak; formative form no better}{NHATS measures them poorly: untested, not refuted}
\hpdbranch{-2.76}{qual}{Invariance fails across sex or age band}{Within-group use only; qualify every subgroup comparison}
\hpdbranch{-3.68}{qual}{Loadings shift with estimator, item or cognition rule}{Structure reported as fragile, or as a range between the two rules}
\hpdbranch{-4.6}{agst}{Domains inseparable; one factor fits as well as three}{Single-factor score only; three-domain theory set aside here}
\node[stepbox, minimum height=5.4cm] (s1) at (\xs,-2.3) {\textbf{1}\\\textbf{Measurement}\\[2pt]Chapter~\ref{ch:results1}\\[4pt]Read first: it decides what step 2 is about};
% ---- Step 2: prognosis
\hpdbranch{-5.85}{kept}{Ranks well; predicted and observed risks agree}{Usable for ranking and at a threshold; net benefit read last}
\hpdbranch{-6.77}{qual}{Death predicted, admission not}{Admission partly measures access; the two outcomes reported apart}
\hpdbranch{-7.69}{qual}{Admission predicted, death not}{Score may track contact with services more than capacity}
\hpdbranch{-8.61}{qual}{Conclusions change on the $[0.05,1]$ floor}{Result carried by the lowest scorers; both rules reported}
\hpdbranch{-9.53}{agst}{Poorly calibrated after internal validation}{Cannot be used at a threshold; counts against the framework}
\node[stepbox, minimum height=4.45cm] (s2) at (\xs,-7.69) {\textbf{2}\\\textbf{Prognosis}\\[2pt]Chapter~\ref{ch:results2}\\[4pt]Discrimination, then calibration, then net benefit};
% ---- Step 3: comparison with frailty
\hpdbranch{-10.78}{kept}{Paired gain survives removal of overlap and changes decisions}{Gain comes from the domains frailty does not cover}
\hpdbranch{-11.7}{agst}{Gain clear of zero, but no decision changes}{Reported as a null for incremental value}
\hpdbranch{-12.62}{agst}{Gain lost when overlapping indicators are removed}{Duplication of frailty information: a null for incremental value}
\node[stepbox, minimum height=2.6cm] (s3) at (\xs,-11.7) {\textbf{3}\\\textbf{Comparison}\\[2pt]with the frailty index, paired};
% ---- Step 4: realisation model
\hpdbranch{-13.87}{kept}{Non-compensatory form beats the additive composite}{Weakest-domain reading gains support in these data}
\hpdbranch{-14.79}{agst}{Indistinguishable from the additive composite, or worse}{Narrowed to a theoretical model; factor score kept}
\node[stepbox, minimum height=1.7cm] (s4) at (\xs,-14.33) {\textbf{4}\\\textbf{Realisation}\\\textbf{model}\\[2pt]Chapter~\ref{chap:crf}};
% Order of reading
\draw[hparrow] (s1) -- (s2);
\draw[hparrow] (s2) -- (s3);
\draw[hparrow] (s3) -- (s4);
% Closing line
\node[hpnote, text width=12.5cm] at (6.5,-15.75) {On every branch the next step is external validation. No branch has yet been observed.};
% Legend
\node[draw, kept, minimum width=0.8cm, minimum height=0.35cm] (k1) at (0.6,-16.45) {};
\node[font=\footnotesize, anchor=west] at (1.1,-16.45) {framework kept};
\node[draw, qual, minimum width=0.8cm, minimum height=0.35cm] (k2) at (4.4,-16.45) {};
\node[font=\footnotesize, anchor=west] at (4.9,-16.45) {kept with a qualification};
\node[draw, agst, minimum width=0.8cm, minimum height=0.35cm] (k3) at (9.0,-16.45) {};
\node[font=\footnotesize, anchor=west] at (9.5,-16.45) {counts against or narrows a claim};
\end{tikzpicture}}

\captionsetup{justification=justified, width=\linewidth}
\caption{Reading plan for the results, fixed before any result is seen. Steps are
read from top to bottom; measurement decides what the later steps are about. Each
row gives a pattern the analysis could return and the reading the chapter assigns
to it. Solid outline, light ground: framework kept. Dashed outline, white ground:
kept with a qualification. Dotted outline, darker ground: counts against the
framework or narrows a claim. No branch has yet been observed.}
\label{fig:disc-decision}
\end{figure}

Any of these results would be reported as evidence against the framework, and the
conceptual claims of this chapter would be narrowed to fit it. They would not be
rescued by appeal to the limitations of the survey, except where the analysis plan
had named that limitation in advance as the reading for that pattern, as it does in
Table~\ref{tab:disc-patterns}.

A finding counts for nothing until it is placed against what the field already
holds. Section~\ref{sec:disc-literature} does that placing, one tradition at a
time, beginning with the theories that decide whether the construct has a subject
at all.

% =====================================================================
\section{Interpretation in the Context of the Literature}
\label{sec:disc-literature}
% =====================================================================

\noindent A new measure is judged against the instruments a reader already
trusts. The question a service asks is not abstract: what would it gain by adding
the $H_{\text{potential}}$ to the frailty index it already scores, and what would
it have to stop doing to make room.

\noindent The findings of this thesis should be interpreted within the broader evolution of health measurement over the past several decades. Traditionally, health has been quantified through disease-specific outcomes, mortality, disability or health-related quality of life. These approaches have transformed clinical research and health policy by providing standardised methods for comparing individuals, interventions and populations. Nevertheless, they remain grounded primarily in descriptions of current health status or accumulated disease burden. The central premise of the present thesis is that these approaches, while individually valuable, do not fully capture an individual's capacity to maintain health into the future. The $H_{\text{potential}}$ was designed to close the biomedical--social gap left by existing measures: QALYs and DALYs summarise burden for economic evaluation, quality-of-life instruments describe current status \citep{whitehead_2010}, and frailty indices quantify physiological reserve \citep{searle_2008}; none integrates biological, social and structural determinants into a single predictive construct grounded in a theory of health \citep{sen_1993,who_2015_ageing,marmot_2005}. Figure~\ref{fig:landscape} positions the construct within this landscape.

\begin{figure}[htbp]
\centering
\begin{tikzpicture}[x=1.25cm, y=1.0cm]
\draw[hparrow] (0,0) -- node[below=2pt, font=\small] {from biomedical only to social and systemic breadth} (9.6,0);
\draw[hparrow] (0,0) -- node[above=2pt, sloped, font=\small] {from current status to future capability} (0,6.2);
\foreach \x/\y/\t/\s in {1.4/1.2/Disease-specific scores/hpbox,
                         4.0/1.6/Quality-of-life instruments/hpbox,
                         2.2/3.6/Frailty indices/hpbox,
                         5.3/3.0/QALY and DALY/hpbox,
                         7.6/5.0/$H_{\text{potential}}$/hpkey}
  \node[\s, minimum height=0.8cm, text width=2.3cm] at (\x,\y) {\t};
\node[hpnote, anchor=north east] at (9.6,-0.75) {Positions are conceptual, not estimated};
\end{tikzpicture}

\captionsetup{justification=justified, width=\linewidth}
\caption{The $H_{\text{potential}}$ relative to established health measures.}
\label{fig:landscape}
\end{figure}

This interpretation therefore considers the contribution of the $H_{\text{potential}}$ within four overlapping bodies of literature: conceptual theories of health, particularly the capability approach and healthy ageing; multidimensional measurement and psychometric modelling; prognostic prediction and risk stratification; and disease-specific measurement, using haemophilia as an exemplar of translational application. Throughout, the emphasis is placed on how the present work extends existing knowledge while acknowledging the substantial contributions of previous research upon which it builds.

The first of the four decides whether the others are worth having. If health is
not a capability, the construct has no subject.

\subsection{Health as Capability Rather than Current Status}
\label{sec:disc-capability}

A measure of current status records what a person does today. It does not
record what the person could still do under other conditions, and that gap is
the subject of this section.

Perhaps the most fundamental contribution of the present work lies in its conceptualisation of health as capability rather than current status. This distinction draws heavily upon Sen's capability approach, in which wellbeing is understood in terms of the opportunities and freedoms available to individuals rather than solely the outcomes they currently experience \citep{sen_1993}. Within health research, this perspective has increasingly influenced thinking about healthy ageing and person-centred care because it shifts attention from the presence of disease towards the capacity to achieve and maintain valued states of health.

The World Health Organization adopted a closely related position in its framework for healthy ageing, defining healthy ageing not simply as the absence of disease but as the process of developing and maintaining the functional ability that enables wellbeing in older age \citep{who_2015_ageing}. Functional ability itself arises through the interaction between intrinsic capacity and the environment within which an individual lives. The $H_{\text{potential}}$ extends this perspective by proposing that health capability can be represented as a measurable latent construct. Rather than treating biological reserve, social context and healthcare systems as independent explanatory variables, the framework integrates these components within a common multidimensional structure. Importantly, this interpretation should not be viewed as contradicting existing definitions of health; instead, it provides an operational framework through which concepts such as intrinsic capacity, resilience and capability may be measured quantitatively.

Capability talk is cheap until something counts it. The field that has come
closest to counting it is frailty research, which reached a workable number by
counting deficits, and the distance between that solution and this one is worth
stating plainly.

\subsection{Relationship to Frailty Theory}
\label{sec:disc-frailty}

Frailty already does part of this job, and does it well enough that any new
measure has to say what it adds. A geriatrician who can already count deficits is
entitled to ask why a second number should be scored on the same patient.

Frailty research provides perhaps the closest conceptual neighbour to the present framework. During the past two decades, frailty has become one of the dominant paradigms within geriatric medicine because it captures the increased vulnerability associated with ageing beyond chronological age alone. Particularly influential has been the deficit accumulation model, in which frailty is quantified through the proportion of accumulated health deficits \citep{rockwood_mitnitski_2007,searle_2008}. The present thesis agrees with frailty theory in recognising that health is a multidimensional phenomenon and that biological reserve is more informative than isolated disease diagnoses when predicting future outcomes. However, the two frameworks differ in their conceptual orientation. Frailty indices principally describe accumulated deterioration: as deficits increase, frailty increases and prognosis worsens. The $H_{\text{potential}}$, by contrast, begins from the opposite perspective, attempting to quantify the remaining capacity available to maintain health rather than the deficits already accumulated. Two individuals with similar levels of frailty may nevertheless differ substantially in healthcare access, social support or treatment opportunities, factors that influence future outcomes but are only indirectly represented within conventional frailty indices.

\medskip
\noindent\emph{[Placeholder: discuss how the observed psychometric structure and comparative performance relate to frailty indices. If the $H_{\text{potential}}$ demonstrated superior discrimination or incremental prognostic value, interpret this as evidence that incorporating social and systemic determinants captures information beyond physiological deficits alone. If performance was comparable, discuss the conceptual implications of obtaining similar predictive performance using a theoretically broader construct.]}

\begin{reading}[What comparable discrimination would mean]
Take two people with the same frailty index. If the analysis returns a
$C$-statistic for the $H_{\text{potential}}$ close to the frailty index and no
better, the reading is not that the extra domains are empty. It is that they
change the ranking too seldom to move an average taken over the whole cohort.
The people in whom they should change it, such as two people at Profile~A who
differ only in social enablement $S$, may be a small share of the sample and are
outweighed by the rest.

If the analysis instead returns a clear gain, what it buys is a reordering.
People at Profile~A move up or down the list of those a service visits first; the number of visits
does not change. That is all a gain in discrimination ever buys, and it buys
nothing unless somebody acts on the order.
\end{reading}

One risk in this comparison is easy to miss. The biomedical domain and the frailty
index draw on some of the same kinds of information, function and physical
performance among them. A gain for the $H_{\text{potential}}$ over the frailty
index could therefore come from a second, slightly differently weighted, measure of
the same deficits, and not from the social and systemic domains at all. The design
guards against this in two ways. The comparison is paired, so both models are
refitted on the same resamples and only their difference is judged. And one
sensitivity analysis removes the indicators that overlap most closely with the
frailty measure; the gain that survives that removal is the gain the broader
domains contribute (Section~\ref{sec:disc-read-prog}). Without that step a positive
result could not be told apart from duplication, and the argument of this section
would rest on a number that does not measure what it is said to measure.

Frailty counts what has been lost. The other established comparator prices what
remains, in money, and it is a different kind of instrument again.

\subsection{Relationship to QALYs and Health Utility Measures}
\label{sec:disc-qaly}

A health utility measure asks what a year of life is worth. The question here
is whether function will be kept over five years. Both are fair questions, and
an instrument built for one answers the other badly.

Quality-adjusted life years (QALYs) have become the dominant metric for economic evaluation within healthcare because they combine survival and health-related quality of life into a single measure of health utility. Similarly, disability-adjusted life years (DALYs) quantify population disease burden. Together these measures have transformed health technology assessment, comparative effectiveness research and resource allocation. Despite their widespread use, QALYs and DALYs were never intended to measure latent health capacity. Their principal objective is normative evaluation of healthcare interventions rather than prediction of individual future outcomes. Utility scores describe current preferences or health states, whereas the $H_{\text{potential}}$ seeks to estimate the capacity from which future health states emerge. Consequently, the two approaches should be viewed as complementary rather than competing representations of health.

\medskip
\noindent\emph{[Placeholder: interpret the comparison between the $H_{\text{potential}}$ and the QALY proxy. Discuss whether the results suggest that prediction of future health requires information extending beyond health-related quality of life, or alternatively whether comparable performance indicates substantial overlap between the constructs.]}

\begin{reading}[Discriminant validity, tested on Profile~B]
Profile~B with structure favourable and with structure unfavourable share the
same biomedical potential and social enablement. One coordinate moves, from
$(0.30,0.70,0.85)$ to $(0.30,0.70,0.20)$, and the $H_{\text{potential}}$ falls
from $0.231$ to $0.112$. A utility score built from symptoms and function on the
day of the interview would move far less, because $P$ and $S$ are unchanged.

If the analysis returns a high correlation between the $H_{\text{potential}}$ and
the QALY proxy, that is evidence the new instrument is measuring the old thing
under a new name. A moderate correlation, with the disagreement concentrated in
people whose circumstances and physiology point in opposite directions, is the
pattern the framework predicts. Where the disagreement sits tells more than how
large it is.
\end{reading}

The comparison has a weaker side, and it should be stated before any result
arrives. NHATS does not administer the EQ-5D. The utility comparator assigns EQ-5D
levels from NHATS items by a crosswalk written for this thesis and values them with
the UK time trade-off tariff, whose coefficients were checked one by one against
their source (Appendix~\ref{app:datachecks}). A correct tariff applied to
crosswalked levels does not make a validated utility. Whatever the analysis returns,
it compares the $H_{\text{potential}}$ with a proxy for a utility score, and a gain
over the proxy is a gain over that proxy, not over the EQ-5D as it would be
administered. The comparator is named a QALY proxy throughout for that reason, and
the reading above should be held to the same limit.

What separates the two instruments is mostly the third domain, the part of a
person's prospects that lies outside the body and outside their stated
preferences. That domain has a literature of its own.

\subsection{Social Determinants of Health}
\label{sec:disc-sdh}

Adjusting for poverty and treating poverty as part of a person's health are
different acts with different consequences. The first takes it out of the score
and puts it in the model; the second leaves it in the score.

One of the distinguishing features of the proposed framework is the explicit integration of social and systemic determinants into the latent construct itself. This position reflects decades of research demonstrating that education, socioeconomic position, housing, employment, social support and access to healthcare shape health outcomes heavily and independently of biological disease \citep{marmot_2005}. Nevertheless, social determinants are commonly incorporated into prediction models only as adjustment covariates. They are treated as variables that modify relationships between biological predictors and health outcomes rather than as constitutive components of health itself. The present thesis adopts a different position by arguing that social and systemic conditions are intrinsic determinants of an individual's capacity to maintain health, and that their influence should therefore be represented within the latent construct rather than added after measurement. The proposed three-domain architecture provides a conceptual bridge between biomedical and social models of health, incorporating them within a unified latent variable rather than viewing them as competing explanations.

\medskip
\noindent\emph{[Placeholder: discuss whether the empirical analyses support the inclusion of Social Modifiers and Systemic Burdens. If these domains exhibited strong factor loadings or improved prognostic performance, relate the findings to the literature on social determinants and healthy ageing. If their contribution proved weaker than anticipated, consider possible explanations, including limitations of available indicators within NHATS.]}

Two tests are easily run together here and should be kept apart. A domain can
matter for outcomes and still load weakly, because the survey measured it through
distant proxies. A domain can load well and still add nothing to prediction, because
what it measures has no bearing on death or admission over five years once the other
domains are known. The first is a question about the indicators; the second is a
question about the construct. A weak result on the contextual domains will be read
by asking which of the two failed, using the rules of
Section~\ref{sec:disc-read-measure}, and not by assuming the more convenient one.

Putting social conditions inside the score has a further consequence. The score
inherits whatever the social and systemic indicators carry, including their
weaknesses. Income, for instance, enters as the average of its logarithm over the
NHATS imputations, which gives one value per person but discards the spread between
imputations, so that standard errors involving that indicator are somewhat too small
(Appendix~\ref{app:datachecks}). An adjustment covariate with the same fault would
distort one coefficient. A domain indicator with that fault passes it into every
person's score. This is a cost of the position taken in this section, and it is
paid in the precision of the contextual domains.

Whether that architecture survives contact with data is a psychometric question
before it is a clinical one, and the psychometric machinery has its own way of
going wrong.

\subsection{Psychometric Measurement of Complex Health Constructs}
\label{sec:disc-psychometrics}

An index whose weights were fitted to death will predict death. What is left when
that shortcut is refused is the question this section answers.

Methodologically, the present thesis contributes to an expanding literature using latent variable models to represent complex multidimensional health constructs. Structural equation modelling has become increasingly important in health research because many clinically relevant concepts, including resilience, quality of life, intrinsic capacity \citep{cesari_2018} and wellbeing, cannot be observed directly but must instead be inferred from multiple imperfect indicators \citep{kline_2015,brown_2015}. The principal methodological contribution of the present work lies in its clear separation of measurement from prediction. Many composite indices derive weights directly from their association with clinical outcomes, creating a form of circular reasoning in which predictors are optimised for the outcomes they subsequently claim to predict. In contrast, the $H_{\text{potential}}$ was first estimated as a latent construct using confirmatory factor analysis (a method that tests whether the indicators vary together in the pattern a stated structure predicts) before any prognostic analyses were undertaken, so that prediction evaluates the independent usefulness of the construct rather than contributing to its definition. This design is formalised below. In plain terms, the derivation shows that the score is fixed from the indicators alone, before any death or admission is consulted, in the way an assay is calibrated against reference standards and not against the diagnoses it will later help to make. It then states the limit of that guarantee: a score built without the outcome can still be tuned to the sample, and that fault needs its own remedy.

\begin{derivation}[Why the two-stage design avoids outcome-driven circularity]
\label{der:noncircularity}
Let $\widehat H_i=m(\mathbf x_i;\widehat{\boldsymbol\theta}_m)$, where the
measurement criterion uses the indicators $\mathbf x$ but not outcome $Y$.
Prediction then estimates $p(Y_i\mid\widehat H_i,\mathbf z_i)$. Thus the
construct weights are not chosen to maximise apparent association with $Y$.
In contrast, a supervised index $I_i=\mathbf w^{\mathsf T}\mathbf x_i$ with
$\widehat{\mathbf w}=\arg\max_{\mathbf w}\ell_Y(\mathbf w)$ reuses outcome noise
when selecting and evaluating its weights. Separation prevents outcome-driven
definition, but does not eliminate overfitting: indicator selection, SEM, scoring,
and prognosis are sample-dependent and must all be repeated within validation.
\end{derivation}

\begin{reading}[What the separation buys, and what it leaves open]
The derivation makes one guarantee. Whatever the score turns out to predict, it
was not built to predict it. Each respondent's three domain scores would be the
same numbers if every death date in NHATS were shuffled, and that is a claim which
can be executed rather than argued about.

What it does not buy is protection from optimism. The indicators were chosen,
the factor structures were compared, the scores were computed and the survival
model was fitted in one sample, and every one of those steps can be tuned to
that sample without ever looking at the outcome. This is why the whole
sequence, imputation included, is repeated inside each bootstrap replicate and
not just the last step of it. A design can be free of circularity and still be
over-fitted, and the two faults need separate remedies.
\end{reading}

\medskip
\noindent\emph{[Placeholder: discuss how the observed factor structure, reliability estimates and measurement properties compare with previous latent constructs reported in the literature. Consider whether the three-domain model proved simpler, more interpretable or more stable than alternative approaches.]}

Refusing the shortcut has a price. The construct reaches the prognostic stage
untuned for it, and must then be judged on terms it was never fitted to meet.

\subsection{Interpretation of Prognostic Performance}
\label{sec:disc-prognostic}

A p-value cannot tell a service how many home visits to fund. Two further terms
carry most of the weight in what follows, and both have clinical readings. Net
benefit is the trade a clinician makes between finding a case and over-treating:
it counts the true cases the model finds and charges it for the false ones, at
the exchange rate implied by the risk at which that clinician would act
\citep{vickers_elkin_2006}. Internal validation asks how much of the reported
performance is real and how much is the model marking its own homework; the
correction it applies, shrinkage, is regression to the mean under another name,
since an estimate singled out because it looked extreme in this sample will look
less extreme in the next \citep{harrell_1996,steyerberg_2019}.

The prognostic analyses presented in Chapter~\ref{ch:results2} represent the first empirical evaluation of whether the latent construct captures clinically meaningful information regarding future health outcomes. Modern prediction research emphasises that clinically useful models require more than statistically significant associations: discrimination, calibration, internal validity and clinical utility are now regarded as equally important components of model evaluation \citep{steyerberg_2019,vancalster_2019}. The present thesis therefore adopted this broader evaluation framework rather than relying exclusively upon regression coefficients or hazard ratios. Prediction models should be judged according to whether they improve decision making, and not on whether their predictors achieve statistical significance, and the inclusion of decision-curve analysis and comparison with established measures positions the present work within contemporary standards for prediction model evaluation. The notion of incremental value used in that comparison is defined below. The definition says two things in symbols. The new score is judged by the change in performance when it is added to the model already in use, not by its performance alone. And because both models are fitted to the same people, the change must be estimated on the same resamples, as a method-comparison study runs both assays on the same specimens so that differences between specimens cancel and only the difference between methods remains.

\begin{derivation}[Incremental prognostic value over a baseline model]
\label{der:incremental}
For baseline predictor $\eta_{0i}=\boldsymbol\gamma^{\mathsf T}\mathbf z_i$ and
extended predictor $\eta_{1i}=\eta_{0i}+\beta\widehat H_i$, define
$\Delta M=M(\eta_1)-M(\eta_0)$ for a performance functional $M$. Examples are
$\Delta C=C_1-C_0$, $\Delta\mathrm{AIC}=\mathrm{AIC}_1-\mathrm{AIC}_0$, and
$\Delta\mathrm{NB}(p_t)=\mathrm{NB}_1(p_t)-\mathrm{NB}_0(p_t)$. Because both
models use the same participants, their estimates are dependent. Paired bootstrap \citep{efron_tibshirani_1993,steyerberg_2019}
replicates must fit both models and retain the within-replicate difference.
Incremental value requires improved decision-relevant performance with acceptable
calibration; a significant $\beta$ alone is insufficient.
\end{derivation}

\noindent For a clinician the last line of the derivation is the one that
matters. A coefficient that is statistically significant shows that the score is
associated with the outcome once the frailty index is known. It does not show
that anyone would be managed differently, and only the second would justify
asking a service to collect the extra information.

\medskip
\noindent\emph{[Placeholder: interpret the observed discrimination, calibration, comparative performance and clinical utility. Discuss how the results compare with published prognostic models in ageing populations and whether the observed performance supports further development of the $H_{\text{potential}}$.]}

\begin{reading}[Net benefit at a fixed threshold]
Decision-curve analysis asks at what predicted five-year risk an intervention,
such as a medication review with a home assessment, would be offered. Suppose
the threshold is 20 per cent. Net benefit at that threshold states the value of the model in one unit,
cases found per hundred people assessed, after charging it for everyone it sends
for assessment who would never have had the event.

If the analysis returns a net benefit for the $H_{\text{potential}}$ above the
frailty index at 20 per cent, the report should be arithmetic: per hundred older
adults assessed, so many more people at the risk of Profile~A are identified for
the same amount of over-assessment. If the two curves cross, the report names the range of
thresholds over which each model wins, because a service that acts at 10 per cent
and a service that acts at 30 per cent are not being told the same thing.
\end{reading}

Two hazard ratios put next to each other can be read in two directions, and
the directions come apart. The next example sets out both, using figures the
analysis has not yet produced.

\begin{workedexample}[disc:incremental]{Reading an incremental hazard ratio against its comparator}
The comparison in Derivation~\ref{der:incremental} is between a baseline model
carrying a frailty index and an extended model carrying the
$H_{\text{potential}}$ as well. Suppose five-year all-cause mortality at the
cohort mean is $20$ per cent, and consider a person one standard deviation into
the worse tail, roughly the position of Profile~A. All figures below are
supposed, not observed.

Case one. If the analysis returns a hazard ratio of $1.35$ per standard
deviation for the frailty index and $1.40$ for the $H_{\text{potential}}$, the
predicted five-year risks for that person are $26.0$ and $26.8$ per cent. The
difference is $0.8$ percentage points, one person in $122$. In a cohort of
several thousand a paired bootstrap could easily put the difference in
discrimination clear of zero, so the construct would be statistically
distinguishable from its comparator and would still change no decision at any
threshold in use. That result should be reported as a null for
incremental value, whatever its confidence interval says.

Case two. If instead the frailty index returns $1.30$ and the
$H_{\text{potential}}$ returns $1.90$, the same person's predicted risks are
$25.2$ and $34.6$ per cent, a gap of $9.4$ points, one person in $11$. That
difference would reorganise a case list. If the paired interval for the
difference spans zero, the study has not shown it, and the correct report is
that a difference of decision-changing size has not been ruled out and has not
been established either. The remedy is a larger sample, not a firmer adjective.

\begin{wecode}[Python]
from hpv.prognostic import hazard_ratio_to_risk, number_needed

base = 0.20
for comparator, construct in ((1.35, 1.40), (1.30, 1.90)):
    r0 = hazard_ratio_to_risk(base, comparator)
    r1 = hazard_ratio_to_risk(base, construct)
    print(f"HR {comparator} vs {construct}: risk {r0:.3f} vs {r1:.3f}  "
          f"difference {r1 - r0:.3f}  one in {number_needed(r0, r1):.0f}")
\end{wecode}

\weanswer{Risks of $0.260$ against $0.268$ in case one, a difference of $0.008$
or one person in $122$; and $0.252$ against $0.346$ in case two, a difference
of $0.094$ or one person in $11$. Statistical distinguishability and clinical
worth are separate findings and can point opposite ways. Asserted by
\texttt{test\_worked\_example\_disc\_incremental\_value\_two\_cases}.}
\end{workedexample}

Performance in a development sample is a claim about one population. The
framework was built to travel, and the hardest place it has to travel to is a
haemophilia clinic.

\subsection{Translation to Haemophilia}
\label{sec:disc-translation}

Nothing in a cohort of community-dwelling Americans over 65 was collected with
haemophilia in mind. Profile~B, with low biomedical potential and a structural
domain that differs between health systems, is therefore a fair test of whether
the construct is general or only well fitted.

The final contribution of the thesis concerns translation beyond the developmental population. Rather than constructing a disease-specific model from first principles, the present work proposed a methodology in which the latent construct remains stable while observable indicators are adapted to reflect disease-specific biology and healthcare systems. This strategy aligns with current recommendations regarding transportability and external validation of prediction models, recognising that theoretical constructs should remain stable while empirical parameters may require recalibration in new populations \citep{moons_2019,collins_2024}. Importantly, the haemophilia application should not be interpreted as evidence that the framework has already been validated within that population; rather, it demonstrates the conceptual feasibility of translation and establishes a structured pathway for future empirical investigation.

Translation raises a second question that the transportability literature does not
answer for clinical models alone. Even if the construct and its recalibrated
parameters carry over to a haemophilia clinic, the decisions built on it need not.
Chapter~\ref{chap:crf} proves that equal capability gains and equal health gains in
two health systems do not imply equal policy rankings, because prices, thresholds
for opportunity cost and budget ceilings differ between them
(Theorem~\ref{thm:nontransport}). The point bears directly on haemophilia, where
the cost of factor concentrates and the organisation of care differ widely between
systems. A score that transports well would still leave each system to decide what
to do with it under its own prices and budget. The translational pathway of
Chapter~\ref{ch:translational} therefore ends at a validated measure, not at a
treatment recommendation, and it is right to stop there.

That is the fourth and most modest of the comparisons. What the four amount to
together still has to be said, because read side by side they can be mistaken for
a larger claim than the work supports.

\subsection{Overall Interpretation}
\label{sec:disc-overall}

Set out one after another, the four comparisons could be read as a claim to have
replaced everything they compare with. They are not that, and the point is worth
closing off before the design is examined.

Taken together, the literature suggests that health measurement has progressed through successive stages, beginning with disease-specific indicators, expanding to multidimensional quality-of-life instruments, incorporating frailty and healthy ageing frameworks, and increasingly recognising the role of social determinants and healthcare systems. The present thesis builds upon each of these developments rather than replacing them. Its principal contribution is to integrate these previously separate traditions within a single latent measurement framework that is theoretically grounded, psychometrically structured and designed for prognostic evaluation. Whether the empirical findings ultimately demonstrate superiority, non-inferiority or complementarity relative to existing measures, the broader conceptual contribution remains unchanged.

That last claim needs a limit, or it would make the conceptual contribution
immune to evidence. The argument that health is a capability with three sources
does not depend on any one cohort. The claim that this capability can be measured
does. If the analysis returns one of the patterns that Section~\ref{sec:disc-read-against}
counts against the framework, the conceptual argument survives as an argument, but
the thesis can no longer say it has shown the construct to be measurable, and it
will not say so. What remains unchanged under every result is the claim about how
health should be conceived. What changes with the results is how much of that claim
has been brought within reach of data.

A conceptual claim is only as good as the design that produced it. What that
design earns, and what it cannot earn, is the subject of
Section~\ref{sec:disc-strengths-limitations}.

% =====================================================================
\section{Strengths and Limitations}
\label{sec:disc-strengths-limitations}
% =====================================================================

\noindent Most of what is right about this study and most of what is wrong with
it come from the same few decisions. Splitting them into two lists is a
convenience of exposition, and where a strength and a limitation are two readings
of one choice, that is said.

\noindent The study joins conceptual theory, psychometric measurement, survival analysis and translational planning in one design, and that breadth brings both strengths and constraints. Many decisions that strengthen one part of the study constrain another. Developing the measurement model within a single well-characterised cohort, for example, improves internal consistency and methodological control but limits immediate generalisability; separating measurement from prediction reduces circularity but introduces additional analytical complexity.

\subsection{Strengths}
\label{sec:disc-strengths}

A strength is only a strength against the alternative a competent researcher
might have chosen instead. Each item below names the alternative it turned down.

\subsubsection{A theory-grounded conceptual framework}

The alternative was to let a regression choose the domains.

The $H_{\text{potential}}$ was built from an explicit theory, not found by statistical optimisation. Many composite indices combine whatever variables are associated with outcomes; they often predict well, but their meaning is unclear because the statistical model defines the construct. Here the three domains came from a synthesis of the capability approach, healthy ageing, frailty science and the social determinants of health, and statistical modelling was used to test whether that structure holds, not to create it. Each domain therefore has a meaning fixed before the data, which supports construct validity and makes the results easier to interpret.

\subsubsection{Separation of measurement and prediction}

The cheapest way to make a score look prognostic is to build it out of the
outcome.

Composite indices are often built by choosing variables or weights against the outcomes they later predict. That maximises apparent performance in the development data but makes the outcome part of the predictor's definition. This study instead estimated the latent construct by structural equation modelling, without reference to any outcome, and only then entered the factor scores into separate survival models. The prognostic analyses therefore test a finished construct instead of helping to build it. This lowers the risk of overstating performance, lets measurement and prognostic properties be judged separately, and eases translation, since the measurement model can be adapted to a new population without changing the concept.

\subsubsection{Integration of biological, social and systemic determinants}

Standard practice enters social factors as covariates, where they adjust the
estimate and never reach the patient's score.

Existing measures usually centre on one aspect of health, such as disease burden, disability, quality of life or physiological reserve, and leave social and healthcare factors to be adjusted for separately. The $H_{\text{potential}}$ treats them as parts of health, in line with the view that outcomes arise from the interaction of biological, behavioural and environmental processes. The structure suits ageing populations, in whom access to care, social support, income and multimorbidity interact over the long term.

\subsubsection{Use of a large longitudinal cohort}

A three-domain latent model needs more indicators than most cohorts trouble to
collect.

The analyses use the National Health and Aging Trends Study, one of the most detailed longitudinal studies of older adults available. Its range of indicators allows a multidimensional latent model that a narrower dataset could not support; its prospective follow-up places measurement of the predictor before the outcomes; and the diversity of its sample steadies the estimation of the measurement model. The cohort is limited to older adults in the United States, but its size, quality and design give the development phase a sound base.

\subsubsection{Prognostic performance judged on more than one axis}

A model can rank patients well and still be worthless at the risk where the
decision is actually made.

Another strength lies in the approach adopted for evaluating prognostic performance. Rather than relying exclusively upon measures of statistical association, the study assessed discrimination, calibration, comparative performance, decision-curve analysis and optimism-corrected internal validation. Comparison with established measures, including frailty indices, comorbidity measures and QALY proxies, places the proposed framework within the context of existing clinical practice, while assessment of clinical utility through decision-curve analysis extends evaluation beyond statistical accuracy to consideration of potential clinical relevance.

\subsubsection{Transparent and reproducible methodology, and a defined translational path}

Analysis plans written after the analysis tend to describe the analysis that
worked.

The study also benefits from a transparent methodological framework in which model specification, comparison strategies, internal validation and handling of missing data were defined prospectively, following recognised standards for prediction models and measurement properties \citep{collins_2024,moons_2019}. Finally, the thesis extends beyond measurement and prediction by proposing a structured translational pathway for disease-specific adaptation. The proposed sequence of indicator mapping, expert consultation, psychometric re-estimation, recalibration and external validation provides a scientifically rigorous methodology for adapting the framework to haemophilia and, potentially, other chronic diseases.

\subsubsection{Data checked against their source documents}

The usual alternative is to trust the variable names in a data release and the
labels attached to them.

Every variable name, file and response code used to build the analytic file was
checked against the NHATS crosswalk, instrument and User Guide before any analysis
(Appendix~\ref{app:datachecks}). The checks found faults that would have changed
the results had they gone unnoticed. Activity codes had been read in reverse, which
would have counted people who manage an activity alone without difficulty as
limited in it, and so inverted part of the function indicator. Education was
missing for every continuing member of the cohort, because NHATS asks it only at a
person's first interview, which would have left one of the four systemic indicators
to be imputed for all of them. Deaths with no proxy interview had been left out of
the mortality outcome and would have been treated as losses to follow-up. None of
these faults would have stopped the analysis from running, and none would have
looked wrong in its output. That is what makes them dangerous, and it is why the
checks were made against documents and not against the plausibility of the results.

The checks are also enforced, not merely recorded. Before the analytic file is
built, the pipeline confirms that every declared variable is present and stops if
one is absent, and the test suite holds a test for each correction, so that a later
change to the code which brought back one of the errors would fail before it reached
the results. These are separate layers of defence, each able to catch what another
misses, in the sense that \citet{reason_1990} gives to defences against error: a
document check catches a wrong name, the pre-build check catches a release in which
a name has changed, and the tests catch an error reintroduced by an edit. A
strength of this kind is invisible when it works, which is a reason to state it.

\subsubsection{Sensitivity analyses fixed in advance}

A sensitivity analysis chosen after the main result is known tends to be the one
that leaves it standing.

The analysis plan names six objections and the analysis that answers each:
overlap with the frailty comparator, alternative factor structures, formative or
hybrid specifications, informative missingness and attrition, the definition of
hospitalisation, and the floor of the domain scale. Fixing these in advance does
two things. It stops the set of checks from being chosen to suit the result, and it
lets each check be read by a rule set before the result, as
Section~\ref{sec:disc-reading} does. The cost is that some checks will be run that
turn out to matter little, which is a small price for checks that cannot be
selected after the fact.

In brief, the strengths that hold regardless of the direction of the results are as follows.

\begin{itemize}
  \item \textbf{Theory-grounded construct.} The three domains are anchored in the capability approach, healthy ageing, frailty science, and the social determinants of health, rather than assembled for statistical convenience.
  \item \textbf{Separation of measurement and prediction.} Deriving the construct by structural equation modelling before, and independently of, the prognostic analysis reduces the circularity that affects composite indices whose weights are fitted to the outcomes they later predict.
  \item \textbf{Integration of social and structural determinants.} The framework treats social and systemic factors as intrinsic components of the health construct, not as adjustment covariates added after the fact.
  \item \textbf{Standardised, transparent reporting.} The study follows recognised standards for prediction models and measurement properties \citep{collins_2024,moons_2019}, with a pre-specified analysis plan and internal validation.
  \item \textbf{A defined translational path.} The haemophilia mapping and external-validation roadmap of Chapter~\ref{ch:translational} give a concrete route from a disease-independent metric to disease-specific use.
  \item \textbf{Data checked and the checks enforced.} Names and codes were checked against the NHATS documents, the faults found were corrected, and a test for each correction runs on every change to the code (Appendix~\ref{app:datachecks}).
  \item \textbf{Sensitivity analyses fixed in advance.} Six objections and the analysis answering each were named before any outcome was used, and each has a reading fixed in Section~\ref{sec:disc-reading}.
\end{itemize}

Each of these decisions bought something and cost something. The costs are set
out next, matched where possible to the strength that produced them.

\subsection{Limitations}
\label{sec:disc-limitations}

The limitations worth stating are those that would change a reader's decision
about whether to use the measure, and they are given in that order.

Notwithstanding these strengths, several limitations should be acknowledged. These limitations do not invalidate the findings but define the scope within which they should be interpreted and identify priorities for future research.

\subsubsection{Development within a single cohort}

Everything the measure knows, it learned from older Americans living in their own
homes.

The most important limitation is that development and internal validation were undertaken within a single cohort. Although NHATS provides an excellent development dataset, participants are community-dwelling adults aged 65 years and older residing in the United States. Consequently, the latent structure and prognostic relationships identified within this study may not generalise directly to younger populations, institutionalised individuals or healthcare systems operating under different organisational, cultural or economic conditions. External validation is therefore essential before broader application of the framework can be recommended.

\subsubsection{Internal rather than external validation}

Bootstrap correction answers a narrower question than it is usually asked to
answer.

In laboratory terms, bootstrap internal validation resembles repeating a validation run many times on material resampled from one batch. It shows how stable the performance is within that batch. External validation corresponds to an inter-laboratory comparison, in which the method meets specimens, staff and conditions it was not developed on.

Related to the previous limitation, the study employed bootstrap internal validation rather than independent external validation. Bootstrap procedures estimate optimism and expected performance within the same underlying population dependably, but they cannot determine whether the model remains valid when applied to fundamentally different patient groups. Optimism-corrected performance estimates should therefore be interpreted as reflecting expected reproducibility within similar populations rather than universal transportability.

\subsubsection{Observational study design}

Nothing in this design can say what would happen if a person's circumstances were
changed, only what follows when they differ.

The prognostic analyses are based upon observational longitudinal data, and the reported associations should not be interpreted as evidence of causal relationships. Although adjustment for major demographic variables reduces confounding, residual confounding arising from unmeasured biological, behavioural or environmental factors cannot be excluded. The $H_{\text{potential}}$ is intended primarily as a prognostic rather than a causal construct, and hazard ratios should not be read as estimates of causal effects.

\subsubsection{Measurement error carried forward}

A factor score is an estimate, and the prognostic model treats it as a reading.
The consequence has a name in each language: attenuation bias to a statistician,
regression dilution to an epidemiologist. Both mean the same thing here, that the
association between the $H_{\text{potential}}$ and mortality is more likely to be
understated than overstated. The direction is worth holding onto, because error
of this kind flattens estimates towards no effect: a positive finding is not an
artefact of it, while a null finding may well be. The laboratory version is
familiar: an imprecise assay correlates less strongly with a clinical outcome
than the analyte itself does, because random error in the reading dilutes the
association.

The latent construct is estimated indirectly through observed indicators rather than measured directly, so estimated factor scores inevitably contain measurement error. Although structural equation modelling explicitly accounts for measurement error during estimation, and the bootstrap validation repeated the complete analytical pipeline, uncertainty associated with latent variable estimation cannot be eliminated entirely. This limitation is inherent to all latent variable models and emphasises the importance of continued psychometric evaluation during future external validation.

\subsubsection{Dependence upon available indicators}

The theory asked for constructs the survey never collected.

The operationalisation of the latent domains was necessarily constrained by the variables available within NHATS. Although the dataset contains extensive information, some theoretically relevant constructs could only be represented indirectly or were unavailable entirely, which may affect the content validity of a domain. The empirical measurement model should therefore be interpreted as one operationalisation of the underlying conceptual framework rather than its only possible representation.

\subsubsection{Assumptions regarding missing data}

Multiple imputation is candid about what it assumes and cannot check the
assumption from the data at hand.

Missing observations were addressed using multiple imputation by chained equations \citep{vanbuuren_2011} under the assumption that data were missing at random. While this represents current best practice and is generally preferable to complete-case analysis, it depends upon assumptions regarding the mechanism underlying missingness. If data were systematically missing in ways not explained by observed variables, some degree of residual bias may remain, and sensitivity analyses under alternative missing-data mechanisms would provide a valuable extension.

\subsubsection{Who is in the analytic sample}

The sample is narrower than ``older Americans'', and in a way that matters for
use.

The analytic sample is restricted to people with a Round~5 sample-person
interview. The restriction is needed: people who died before Round~5 had been given
a positive analytic weight in error, and they have no Round~5 interview from which
baseline indicators could be built (Appendix~\ref{app:datachecks}). It also removes
people represented only by a facility questionnaire. The measure is therefore
developed among people living in the community and well enough, or with a proxy
willing enough, to complete an interview at baseline. People in residential care,
among whom low capacity is likely to be common, are absent by design. A score
developed without them cannot be assumed to rank them correctly, and any use in
residential care would need its own validation.

\subsubsection{Shortcomings recorded in the data checks}

The checks of Appendix~\ref{app:datachecks} settled most questions about names and
codes. Six were not settled, or cannot be settled from documents, and each bears
on how the results should be weighed.

\emph{Timing of events.} A death with a proxy interview is dated to the month from
the interval the proxy reports. A death without one is known only to lie between
the last month the person was seen alive and the month the death was recorded, and
is placed at the midpoint. A first hospital stay is likewise placed midway between
two interviews. Errors of this kind are of the order of a few months and fall in
both directions. They matter little for risk at five years and more for risk at one
year, so any result reported at short horizons carries more of this error than the
five-year results do. The exact date of death would remove it for deaths, but it is
a restricted variable and was not used.

\emph{Proxy respondents and cognition.} Social network size is not asked when a
proxy answers, and the cognitive tests are not given when the proxy says the person
cannot be asked. These values are missing for reasons tied to the person's health,
which is the case the missing-at-random assumption handles worst. The imputation
model adds whether a proxy answered, whether the tests could not be given, whether
dementia has been reported, and the AD8 proxy screen. The AD8 is the useful
addition, since proxies answer it for people who were tested and people who were
not, so the model can learn how test scores fall as the AD8 rises and carry that
relation across. But the untested sit at the severe end of the AD8, beyond most of
the people from whom the relation was learned, and the imputation is an
extrapolation. It is likely to overstate their cognition, while the lowest-score
sensitivity rule is likely to understate it for some. The truth lies between the two
rules, and the biomedical domain for these people is known only within that range.

\emph{Income.} Averaging the logarithm of income over the NHATS imputations
discards the spread between imputations, so standard errors that involve the income
indicator are somewhat too small. The square-root scale used to put couple income
on a per-head basis is a convention, not an estimate.

\emph{Skipped home items.} Reading a skipped floors question as one level is exact
for one-floor homes but an assumption for multi-unit buildings, where a dwelling may
span more than one level. The effect is to understate barriers for a minority of
people in such buildings, and so to overstate their structural freedom slightly.

\emph{First-interview items.} Education and gender for continuing members come from
Round~1, four years before the index assessment. Education is unlikely to change at
these ages over that interval, so the cost is small; a continuing member whose
Round~1 answer is itself missing stays missing.

\emph{Documents and release.} The User Guide consulted is the edition covering
Rounds 1 to 14. If NHATS has reissued a Round~5 file since its first release, names
and codes may differ from those documented at the time. The pre-build check guards
against changed names and codes, but not against a change in what a variable means.

None of these six is large enough on its own to overturn a clear result. Their
effect is to widen the uncertainty around a result in ways the reported intervals do
not fully show, and to make results at short horizons and for proxy-reported people
less firm than results at five years for people who answered for themselves.

\subsubsection{Outcome definitions}

The hospitalisation outcome is defined by what NHATS asks, not by what a clinician
would choose.

NHATS asks whether the person has had an overnight hospital stay since the last
interview, and does not separate planned from unplanned admissions. The operational
endpoint is therefore the first overnight stay of any kind. A planned admission for
a hip replacement and an emergency admission after a fall count alike, although only
the second is what a measure of declining capacity is meant to foresee. The
sensitivity analysis on outcome definitions varies this choice but cannot recover a
distinction the survey did not record. Death after one or more missed interviews is
treated as censoring, not as a competing event, because whether a stay occurred in
the unobserved interval is unknown; this is the cautious choice, and it loses some
information about the people whose follow-up was broken.

\subsubsection{Comparators built for this thesis}

The comparators are not the instruments a clinician would use under those names.

The utility comparator is a crosswalk from NHATS items to EQ-5D levels, valued with
a tariff that was checked against its source, and it is a proxy for a utility score,
not a utility score. The frailty comparator is built from NHATS items, and three of
its variable names had to be corrected during the checks. A gain or loss against
either is a gain or loss against the version built here. If the published
instruments were administered directly, the comparison could come out differently,
and the thesis cannot say in which direction.

\subsubsection{One measurement occasion}

The theory describes a capacity that changes over time. The analysis measures it
once.

The $H_{\text{potential}}$ is estimated from the Round~5 interview alone. A single
measurement can rank people by their capacity at that moment but cannot show how
capacity moves, how fast it is drawn down after an illness, or whether it recovers.
These are the properties a systems reading of the construct cares most about
(Section~\ref{sec:disc-systems}), and none of them is estimated here. The dynamical
extension of Appendix~\ref{app:dynamical} describes how they could be modelled, but
it needs repeated measurement that this thesis does not use.

\subsubsection{The floor of the domain scale}

The rule that places the lowest scorers on each domain at zero is a choice with a
visible effect under the geometric form.

Each domain is mapped onto $[0,1]$ using the 1st and 99th percentiles of the
development split, with values outside clipped. Under the geometric realisation
function a person placed at zero on one domain has a realised score of zero whatever
their other domains are. For that small group the score says nothing about the other
two domains, and two such people with very different biomedical potential are scored
alike. The fixed-floor analysis on $[0.05,1]$ shows whether the conclusions depend
on this, but the primary rule remains the one preregistered, and its effect on the
lowest scorers is a property of the measure that users should know.

\subsubsection{Normative and conceptual choices}

Some of the decisions in this thesis are not empirical questions and cannot be
settled by any amount of data.

Finally, the development of any multidimensional health construct necessarily involves normative decisions. The selection of latent domains, the choice of observable indicators, reverse-coding of Systemic Burdens and standardisation procedures all reflect theoretical judgements regarding how health should be conceptualised and measured. These decisions were made transparently and grounded in established literature, but alternative conceptualisations remain possible, and the framework should be viewed as one theoretically informed representation of latent health capability rather than the only valid approach.

In brief, the limitations that qualify the conclusions are as follows.

\begin{itemize}
  \item \textbf{Development in a single cohort.} The metric is developed and internally validated in the NHATS cohort of community-dwelling United States adults aged 65 and older; generalisability to other settings, age groups, and health systems requires external validation.
  \item \textbf{Internal, not external, validation.} Performance is assessed by bootstrap internal validation; external validation in an independent cohort is a distinct and necessary later phase.
  \item \textbf{Observational prognostic analysis.} The prognostic models are associational; residual confounding cannot be excluded, and the hazard ratios should not be read as causal effects.
  \item \textbf{Measurement error carried forward.} The $H_{\text{potential}}$ enters the prognostic model as an estimated factor score; the two-stage design still carries measurement uncertainty into the prediction stage.
  \item \textbf{Indicator availability.} The domains are operationalised with the indicators NHATS provides, so some theoretically relevant constructs may be measured only in part.
  \item \textbf{Assumptions in missing-data handling.} Multiple imputation assumes data are missing at random; departures from this assumption would bias the estimates.
  \item \textbf{Normative choices.} Reverse-coding, standardisation, and the treatment of the systemic burden domain embed normative choices about how disadvantage should count.
  \item \textbf{Sample restricted to Round~5 respondents.} Decedents before Round~5 and facility-only cases are excluded, so people in residential care are not represented.
  \item \textbf{Shortcomings left by the data checks.} Event times are approximate to a few months, proxy-reported cognition rests on an extrapolation, income standard errors are somewhat too small, and some home barriers are understated (Appendix~\ref{app:datachecks}).
  \item \textbf{Outcome and comparator definitions.} The hospital outcome does not separate planned from unplanned stays, and the utility and frailty comparators are versions built from NHATS items.
  \item \textbf{One measurement occasion and a floor at zero.} Capacity is measured once, so its change over time is not estimated, and the lowest scorers on any domain are placed at zero, where the geometric form ignores their other domains.
\end{itemize}

Held inside those bounds, the work still changes what can be said about measuring
health and about how a service might use such a measure. Those consequences are
taken in turn.

% =====================================================================
\section{Implications}
\label{sec:disc-implications}
% =====================================================================

\noindent An instrument that is not yet validated can still change what people
look for. That is the level at which the following claims are pitched, and they
are kept there on purpose.

\noindent The implications reach beyond a single instrument, but they must stay in proportion to what one development cohort can show. They are taken at three levels: how health is measured, how such measures are built, and how the framework might later be used in clinical practice and health systems.

\subsection{Implications for Health Measurement}
\label{sec:disc-impl-measurement}

The question an instrument asks shapes the answer a service gives. Ask what is
wrong with a person and the reply is a list of diagnoses. Ask what the person
could still do and the reply has to include social and structural conditions.

The main implication is conceptual. The thesis argues that health is a latent capability, a person's capacity to maintain health over time, and is more than a description of current disease, disability or quality of life. Existing measures remain valuable for their own purposes, but they describe realised health, not latent capability. The $H_{\text{potential}}$ sits in a different position: it joins dimensions of health usually treated apart, and does not replace the measures that treat them. If external validation supports it, the approach could help move health measurement from static descriptions of disease towards dynamic accounts of capability, in step with work on healthy ageing, systems medicine and personalised care.

\medskip
\noindent\emph{[Placeholder: interpret how the empirical findings strengthen or modify this conceptual contribution. If the measurement model demonstrated strong psychometric performance and meaningful prognostic value, discuss how the evidence supports conceptualising health as a multidimensional latent capability. If the findings suggest refinement of the framework, discuss how the theoretical model may evolve while preserving its central principles.]}

Changing what is measured helps only if other groups can build measures the same
way, which turns the argument into one about method.

\subsection{Implications for Methodological Research}
\label{sec:disc-impl-method}

The two-stage design is transferable even if this particular construct is not.
Any group building a composite measure faces the same temptation to fit its
weights to the outcome it wants to predict.

The methodological contribution of the thesis extends beyond the proposed construct itself. The research demonstrates a structured strategy for developing multidimensional health measures in which conceptual development, measurement modelling and prognostic evaluation are treated as distinct but interconnected stages. Separating measurement from prediction represents an important methodological principle, and the same approach could be applied to the development of other multidimensional constructs within chronic disease, rehabilitation, healthy ageing and public health. The study also demonstrates the value of combining structural equation modelling with modern prognostic model evaluation, connecting two methodological traditions that have largely developed independently. Finally, the translational framework proposed for haemophilia illustrates a general methodology for adapting latent constructs across clinical contexts by preserving the latent architecture while replacing observable indicators and recalibrating measurement parameters.

Two further practices in this thesis transfer to other work whatever its results.
The first is checking a secondary dataset against its source documents and turning
each correction into a test that runs on every change to the code. Survey releases
carry reversed codes, items asked only at first interview, and status codes that
look alike and mean different things; Appendix~\ref{app:datachecks} records
instances of each in NHATS. An analysis can run without error on a file built with
such faults and return plausible numbers, so the faults are found only by looking
for them. The second is fixing the reading of each result before the result is
known, as Section~\ref{sec:disc-reading} does. Fixing the analysis in advance
stops the analysis from being chosen to suit the result. Fixing the reading in
advance stops the interpretation from being chosen to suit it, and closes the gap
through which a weak result is most easily explained away.

A method is taken up when its effect on practice can be seen.

\subsection{Implications for Clinical Practice}
\label{sec:disc-impl-clinical}

Two patients with identical problem lists and different lives do not have
identical prospects, and no instrument in routine use records the difference.

The study tests no intervention, but it bears on future practice. If externally validated, the $H_{\text{potential}}$ could give clinicians a fuller account of a patient than measures of disease burden or impairment. It could help in risk stratification, by finding people whose prospects differ despite similar clinical pictures. It could prompt assessment read across domains. And because its domains span physiology, social life and access to care, it points towards care shared among clinicians, nurses, physiotherapists, psychologists, social workers and service providers. It is not yet a validated clinical tool. External validation, recalibration and prospective implementation studies must come before routine use.

\medskip
\noindent\emph{[Placeholder: following completion of the empirical analyses, discuss whether the observed discrimination, calibration and clinical utility justify progression towards external validation and prospective implementation studies. Conclusions should remain proportional to the strength of evidence obtained within the developmental cohort.]}

\begin{reading}[What a stratification gain would change]
Suppose a service can visit 100 of its 1000 patients this quarter. It ranks them
by a frailty index and visits the top hundred. Add the $H_{\text{potential}}$,
and if the analysis returns a modest gain in discrimination the practical effect
is that some names on the list change places, perhaps eight or twelve of the
hundred. A person at Profile~A, whose biomedical potential alone would not
place them on the list, may now appear on it because of low values on the two
contextual domains.

The gain is realised only if the visit can act on the reason the person was
flagged, and a clinical visit does not change structural conditions. This is the awkward part of putting
systemic burdens inside a clinical score: the score gets better at finding people
a clinic cannot help by clinical means, which is an argument about who else
should be in the room, not an argument against the score.
\end{reading}

A second condition on clinical use concerns calibration in the place of use. A
model calibrated in NHATS predicts risks at the level of mortality and admission
in that cohort. A clinic whose patients die or are admitted at a different rate
would receive risks that are wrong on average even if their ranking is right. For
ranking alone this does not matter. For any rule that acts at a stated risk, it
does, and the model would need to be recalibrated to local event rates before a
threshold drawn from this thesis could be used. The domain scores also call for
care in how they are shown. A clinician shown a low systemic score is being told
something about the person's circumstances, not about their body, and a report
that gives only the combined score would hide which of the two the clinic is being
asked to respond to.

What holds for one patient in front of a clinician holds again for a population
in front of a commissioner, where the same property of the score becomes
contentious.

\subsection{Implications for Health Policy}
\label{sec:disc-impl-policy}

A measure with disadvantage inside it will, if used to allocate, move resources
towards the disadvantaged. Depending on who is asked, that is the point of the
measure or the objection to it.

The framework also bears on population health and planning. Health systems try to find the people at greatest risk so as to direct preventive resources, often using diagnoses or past use of care alone. The $H_{\text{potential}}$ proposes that broader determinants add to future risk, and if later research confirms this, assessment across domains could support wider prevention. Policy needs strong evidence, however, and internal validation in one cohort is a necessary step, not a sufficient one. The thesis adds to the evidence on which policy could later draw; it does not supply grounds for redesigning a health system now.

Two further cautions follow from the work itself. The first is the
non-transportability result of Chapter~\ref{chap:crf}
(Theorem~\ref{thm:nontransport}): a capability gain that justifies a programme in
one health system need not justify it in another, because costs, thresholds and
budgets differ. Policy evidence built on the $H_{\text{potential}}$ would therefore
have to be rebuilt in each system that uses it, even where the measure itself
transports. The second concerns what happens when a score with structural
disadvantage inside it is used to allocate. Used to direct extra help to people
with low scores, it tends to narrow the gaps it measures. Used to direct scarce
treatment to people judged most likely to benefit, it can widen them, since low
structural freedom lowers the score and a low score then lowers the chance of
receiving the help that might raise it. The score does not choose between these
uses; the rule attached to it does. Section~\ref{sec:disc-systems} sets out why
this is a property of the loop formed by the score and its users, and why it
cannot be fixed inside the measurement model.

\medskip
\noindent\emph{[Placeholder: sharpen the implications once the results are known, keeping claims within what internal validation in a single cohort supports.]}

Settling that objection needs evidence from more than one cohort, and none of the
implications above can be pressed until it exists. But the objection also shows
something the measurement chapters could not: the score and the people who use it
form a system of their own, and some of its effects belong to that system, not to
either part.

% =====================================================================
\section{A Systems View of the Thesis and Its Claims}
\label{sec:disc-systems}
% =====================================================================

\noindent The implications above treat the $H_{\text{potential}}$ as an instrument
and its users as standing outside it. The last of them showed that the separation
does not hold once the score is used to allocate. Chapter~\ref{ch:systems} set out a
systems view of health capacity and of the instrument that measures it. This
section applies that view to the thesis as a whole and asks four questions of it:
what kind of property capacity is, where a change does most, what the score does
once people act on it, and how far claims of this kind can be trusted.

A clinical reader can take the section's conclusions without its algebra. The
three domains combine by multiplication, so a very low value on one of them holds
down the whole score, and help does most where a person is weakest. Once a
service acts on the score, the score and the service form a loop, and the rule
the service uses decides whether gaps between people widen or narrow. None of
this is estimated in the thesis; it is stated so that later studies know what to
measure.

\subsection{Capacity as a Property of the Whole}
\label{sec:disc-sys-emergence}

Haematologists know a quantity of this kind. In haemophilia~B, normal levels of
the other clotting factors do not make up for a lack of factor~IX: the deficient
factor limits the result, and a surplus elsewhere does not offset it. The
realisation function treats the domains of health capacity in a similar, though
milder, way. The parallel concerns the form of the combination only; the domain
levels are not measured concentrations, and the degree to which one domain can
offset another is a modelling choice, not a measured property.

A property of a system is emergent when it belongs to the whole and cannot be
assigned to any part \citep{vonbertalanffy_1968}. Under the geometric realisation
function of Chapter~\ref{chap:crf}, $H=P\,S^{\alpha}R^{1-\alpha}$, health capacity
is a property of this kind. It is a product, not a sum, so no domain carries a fixed
share of it, and the effect of a change in one domain depends on the levels of the
other two. A rise in social enablement adds much to a person with high biomedical
potential and high structural freedom, and little to a person for whom either is
near zero.

This has a direct consequence for how the results are read. A loading, or a hazard
ratio for one domain in a model that holds the others, describes that domain's
association with the outcome on average across the cohort. It is not the domain's
share of anyone's capacity, and a domain with a modest average association may
still be the binding constraint for the people in whom it is lowest. The same point
separates the two scores the thesis carries. The factor score measures what the
indicators have in common. The realisation function is a model of how the domains
combine. If the two differ in what they predict, the difference is information about
combination, not a fault of measurement, and the head-to-head protocol of
Chapter~\ref{chap:crf} is the place where that difference is read.

In the terms of Section~\ref{sec:sys-stocks}, capacity is a stock and realised
health is a flow drawn from it. The interview at Round~5 observes indicators of the
stock at one moment; the outcomes over five years are flows that the stock is meant
to foresee. Two consequences follow. The delay between a change in the stock and a
change in outcomes is built into the design: a measure of capacity is judged by
events years later, and a change made today would show in those events only after a
similar delay. And a single observation of a stock cannot show how fast it is drawn
down or refilled. Resilience in older adults, understood as the ability to resist or
recover from functional decline \citep{whitson_2016}, is a property of that refill,
and it lies outside what one measurement occasion can see
(Section~\ref{sec:sys-resilience}).

\subsection{Where a Change Does Most: the Weakest Domain}
\label{sec:disc-sys-weakest}

This subsection asks where a small improvement would raise the score most, and
answers with the slopes of the realisation function. The answer is a statement
about the model, not about any intervention, and its use for planning is limited
in the ways set out at the end.

The geometric form gives a plain answer to where a small change does most, and the
answer is not the same in relative and absolute terms. In relative terms, a one per
cent rise in $P$ raises $H$ by one per cent, a one per cent rise in $S$ by $\alpha$
per cent, and a one per cent rise in $R$ by $1-\alpha$ per cent, whatever the
person's levels. In absolute terms the marginal gains are $H/P$, $\alpha H/S$ and
$(1-\alpha)H/R$, so a unit of a domain is worth most where that domain is lowest.

Profile~B with an unfavourable structure, $(0.30,0.70,0.20)$, shows the size of
the effect. With $\alpha=0.5$ the gain from a unit of structural freedom, set
against the gain from a unit of social enablement, is $S/R=0.70/0.20=3.5$. A unit
of help spent on the person's structural conditions is worth three and a half times
the same unit spent on their social support, and the ratio grows as $R$ falls.
Under the bottleneck limit of the same family, as $\rho\to-\infty$, only the
smaller of $S$ and $R$ matters, and a gain in the larger is worth nothing at all. The
geometric form sits between this extreme and an additive rule, under which a unit is
worth the same wherever it is spent. The weakest-domain constraint is therefore a
matter of degree, and the degree is set by $\rho$, a parameter the data cannot
estimate and the thesis reports across a range.

Two limits govern how far this ranking can guide action. The first is the one
Chapter~\ref{chap:crf} states: the gradients are model-implied gains, not causal
effects, and acting on them needs separately identified response functions. The
second comes from the feedback set out in Section~\ref{sec:sys-loops}. The gradient
moves one domain and holds the others fixed. If better structural access improves
biomedical status over time, or poor biomedical status erodes social ties, the
domains move together, and the gain from a change in one is larger or smaller than
the gradient says. One measurement occasion cannot estimate those links.

The ranking also says something about the kind of change it describes.
\citet{meadows_1999} orders the places to intervene in a system and puts changes to
numbers and parameters among the weakest, below changes to feedback, to rules and to
goals. A domain level is a parameter in that sense. The more powerful changes lie in
the rules by which services decide who receives what, and those rules are where the
score meets its users (Section~\ref{sec:sys-leverage}). The head-to-head protocol
bears on this too. If the non-compensatory form predicts better than the additive
composite, the weakest-domain reading gains support in these data. If it does not,
the ranking of interventions by the gradient loses its empirical backing, and
policy should not read bottlenecks off it.

\subsection{What the Score Does Once People Act on It}
\label{sec:disc-sys-control}

A score that nobody acts on has no effects beyond its own accuracy. A score that
is acted on becomes part of a control structure, in which controllers take actions
on a process on the basis of a model of that process and receive feedback about the
result. \citet{leveson_2011} treats many failures in such structures as failures of
control: actions taken on a process model that no longer matches the process, or
feedback that is missing, late or misread. Section~\ref{sec:sys-control} sets out
the structure for the $H_{\text{potential}}$. Its controllers would be clinicians
and commissioners; its actions, visits, referrals and allocations; its feedback,
deaths and admissions in the people scored. Four effects follow that no measurement
property of the score can prevent.

The first is delay. The outcomes that show whether the score was used well arrive
over five years. A service that targets the wrong people learns of it late, and by
then it may have acted on several further rounds of scores. Any implementation
would need earlier signals, such as regular checks of whether predicted and
observed event rates still agree, so that the controller is not steering on
information years old.

The second is that acting on the score changes the data from which it would be
updated. If flagged people receive help and the help works, their outcomes improve,
and in later data the score will appear to predict less well than it did. A
recalibration that ignored who had been helped would then weaken the score at the
point where it had been most useful. The process model and the process drift apart
because of the controller's own action, which is the failure Leveson describes.
Updating would need to record who was acted on and why.

The third is the loop described under policy. With structural disadvantage inside
the score, an allocation rule that favours high predicted benefit lowers the help
reaching people with low structural freedom, which can lower their structural
freedom further, which lowers their score at the next round. That is a reinforcing
loop in the sense of Section~\ref{sec:sys-loops}. A rule that directs help to low
scores closes a balancing loop instead. The score is the same in both cases; the
loop is set by the rule. This is why the ranking of places to intervene above put
rules above parameters.

The fourth is pressure on the indicators. The measurement model treats each
indicator as a reflection of the latent domain behind it. Once resources follow the
score, a service has reason to act on what is recorded, a contact logged or a form
completed, and the recorded value can then rise without the capacity behind it
changing. The indicator stops reflecting the domain, and the measurement model that
justified it no longer holds for the people the service has touched. Nothing in the
NHATS analysis can show this, because no one in NHATS was scored or acted on by the
$H_{\text{potential}}$.

These effects are requirements for any implementation, not findings of this thesis.
In the terms of Section~\ref{sec:sys-pipeline}, an implementation would need to
report the domain scores beside the combined score, so that the controller can see
which domain produced a low value; to state the decision rule with the score,
because the rule fixes the direction of the loop; to monitor calibration by domain
and by equity stratum over time; and to record the actions taken, so that later
updating can separate the score's accuracy from the service's effect. None of these
requirements has been tested here.

\subsection{The Limits of These Claims}
\label{sec:disc-sys-limits}

The systems reading adds questions the statistical analysis would not raise. It
does not answer them, and its own claims need limits stated as plainly as those of
the survival models.

The loops, delays and control effects described above are drawn from theory. None
is estimated. A single measurement occasion cannot tell a reinforcing loop between
two domains from a common cause acting on both, and it cannot measure a delay at
all. Every model differs from the system it represents, and making the model more
elaborate does not close that gap \citep{box_1976}; a diagram of loops is a model in
this sense and is held to the same standard. The boundary drawn in
Section~\ref{sec:sys-boundary}, around a person, three domains and the services that
act on them, is a choice. \citet{checkland_1981} treats a system as a way of
organising inquiry that the inquirer chooses, not a thing found in the world, and a
boundary drawn around the household or the neighbourhood would place the binding
constraints and the useful places to act elsewhere. Complex-systems methods have
been proposed for the study of health inequalities \citep{diez_roux_2011}, but they
too depend on data about change over time, which this thesis does not have.

Section~\ref{sec:sys-limits} sets these limits out in general. For this chapter
they mean that every statement in the present section is a hypothesis: for the
dynamical programme of Appendix~\ref{app:dynamical}, which could estimate stocks,
flows and delays from repeated measurement, and for implementation studies, which
could observe the control structure in use. The systems view earns its place by
naming what those studies must measure. It does not license any claim that the
NHATS analysis has not made.

Those studies, and the external validation that must come before them, are the
subject of the next section.

% =====================================================================
\section{Future Work}
\label{sec:disc-future}
% =====================================================================

\noindent The thesis ends with a construct tested once, in one country, against
two outcomes. What follows is what would have to happen before it could be
trusted anywhere else.

\noindent The thesis lays the conceptual and methodological foundations of the $H_{\text{potential}}$ and points to the studies that would refine, extend and test it. Those studies are the next stage of the same programme, not additions to it.

The immediate priority is external validation. Internal validation provides an important assessment of optimism within the developmental cohort, but external validation remains the definitive test of whether the latent construct and its associated prognostic relationships generalise beyond the original study population. Independent validation should therefore be undertaken in cohorts differing in geography, healthcare systems, demographic composition and disease profiles, determining whether recalibration alone is sufficient or whether more substantial modification of the measurement model is required.

The second priority concerns disease-specific translation. This thesis proposed haemophilia as the first translational application because it exemplifies the interaction between biological disease, social context and healthcare systems. The next phase of research should implement the roadmap outlined in Chapter~\ref{ch:translational}, beginning with a South African haemophilia cohort, before extending the approach to other chronic diseases of ageing and multimorbidity, including neurodegenerative disorders, chronic kidney disease and diabetes.

A third direction concerns refinement of the measurement model itself. Future datasets incorporating richer biological, behavioural and environmental information may permit expansion or modification of the indicator set while preserving the underlying latent architecture. Advances in wearable technologies, digital health monitoring and electronic health records may provide new opportunities to measure health capability more continuously than is currently possible using traditional survey-based assessments. Longitudinal modelling represents a related area for development: the present thesis conceptualises the $H_{\text{potential}}$ as a latent construct estimated at a particular point in time, yet health capability is inherently dynamic, and repeated measurement would permit investigation of trajectories of change, transitions between levels of health potential and responses to clinical intervention.

The dynamical framework introduced conceptually in Appendix~\ref{app:dynamical} provides one possible direction for this development. Extending the current latent construct through systems of ordinary and stochastic differential equations offers the possibility of modelling the temporal evolution of health potential under varying biological and environmental conditions, and the haemophilia-specific dynamical extension presented in Appendix~\ref{app:hemo_dynamics} similarly provides a theoretical basis for integrating clotting-factor kinetics, inhibitor development and emerging gene therapies into future mechanistic models. These developments should be regarded as separate research programmes requiring independent theoretical development and empirical validation rather than as components of the present thesis. Longer-term methodological developments may also prove valuable: Bayesian measurement models could incorporate prior biological knowledge into estimation of latent health capability, machine learning approaches may assist with identifying informative indicators while preserving theoretical interpretability, and causal modelling, inverse problems and optimal control theory may eventually permit estimation of intervention strategies capable of maximising future health potential. At present, these directions remain speculative and should be viewed as opportunities for future investigation rather than established extensions of the current framework.

Some of the shortcomings recorded in Appendix~\ref{app:datachecks} can be removed
with access the present work did not have, and doing so would sharpen the next
round of analysis. The exact date of death, held in the restricted NHATS files,
would replace the midpoint placed between two interviews and firm up any result at
short horizons. Carrying the separate NHATS income imputations through the whole
pipeline, in place of their average, would give standard errors that reflect the
uncertainty in income. And a cohort that measures cognition directly in people now
reached only through proxies would replace an extrapolation with an observation.
Each of these is a narrow task, and each removes an uncertainty that the present
results can only bound.

The systems reading of Section~\ref{sec:disc-systems} adds one more direction. Any
implementation study should treat the score and its users as one system and observe
both: which decisions the score informs, under what rule, and with what effect on
the people scored and on the score's later accuracy. Without that, an
implementation can show whether outcomes improved but not whether the score, the
rule or the service produced the change.

In brief, the principal directions for future work are as follows.

\begin{itemize}
  \item \textbf{External validation.} The immediate next phase is external validation of the $H_{\text{potential}}$ in independent cohorts, and recalibration where needed.
  \item \textbf{Disease-specific adaptation.} The haemophilia roadmap of Chapter~\ref{ch:translational} is the first disease-specific application, beginning with a South African haemophilia cohort, and the approach extends to other chronic conditions of ageing and multimorbidity.
  \item \textbf{Dynamical extension.} The static construct can be generalised to a time-varying form. The dynamical ODE--SDE--PDE architecture set out in Appendix~\ref{app:dynamical}, and its haemophilia specialisation in Appendix~\ref{app:hemo_dynamics}, are directions for future modelling, to be developed and validated separately from the measurement-and-prediction work reported here.
  \item \textbf{Longer-term methodological extensions.} Later work could explore inverse modelling and Bayesian calibration to estimate the construct directly from large longitudinal data, and control-theoretic approaches to time-dependent intervention. These are theoretical directions and are not claims of the present thesis.
  \item \textbf{Closing the data gaps.} The restricted date of death, income uncertainty carried through every imputation, and direct measures of cognition in proxy-reported people would remove uncertainties the present results can only bound.
  \item \textbf{Studying the score in use.} Implementation studies should observe the decision rule, the actions taken and their effect on the score's later accuracy, as well as patient outcomes.
\end{itemize}

Those are the directions the work opens. Before the chapter states what it has
settled, the contribution of the thesis as a whole is set out in one place.

% =====================================================================
\section{What This Thesis Contributes}
\label{sec:disc-contributes}
% =====================================================================

\noindent A new assay earns its place in two stages. Its analytical validation,
precision, linearity and agreement with reference material, can be completed at
the bench. Its clinical value can be shown only in patients. The contributions of
this thesis divide the same way. Five stand now, on the specification, the
mathematics, the code and the argument, and can be checked by reading them or
running them. Four depend on the NHATS results and are stated here as what those
results would establish if they fall on the supported branches of
Figure~\ref{fig:disc-decision}. The novelty table in A Guide for the Reader sets
the same claims out chapter by chapter, with what existed before each.

\paragraph{Contributions that stand now.}
The first is the specification of the construct. Capability theory, frailty
indices and health utilities existed before this work; none defines health
capacity as one latent quantity with three named domains, Biomedical Potential,
Social Modifiers and Systemic Burdens, each tied to stated survey indicators and
estimated without reference to any outcome (Chapters~\ref{ch:theory}
and~\ref{ch:methods}). The specification can be checked by reading it, and it fixes
what the measurement model is required to show.

The second is the realisation model and its proofs (Chapter~\ref{chap:crf}).
Constant-elasticity functions and geometric means are established tools, and no claim is
made for them. What the thesis adds is their use as a stated rule for combining
the three domains, derived from axioms, with a theorem that characterises the
primary form (Theorem~\ref{thm:crf_characterisation}), an exact decomposition of a
person's capability gap by domain, a budget-constrained policy gradient, a proof
that policy rankings need not carry over between health systems
(Theorem~\ref{thm:nontransport}), and a head-to-head test fixed in advance that
could reject the model (Section~\ref{sec:E-falsification}). The proofs can be
checked line by line, and each stated property and worked example is asserted by
a test in the verification suite (Appendix~\ref{app:verification}).

The third is the verified analysis pipeline and the checks of the data. Every
variable name, file and response code was checked against the NHATS documents,
and the checks found faults that would have changed the results without any sign
in the output: reversed activity codes, education missing for every continuing
member of the cohort, and deaths without a proxy interview left out of the
mortality outcome (Appendix~\ref{app:datachecks}). Each correction is held by a
test that runs on every change to the code. In laboratory terms this is quality
control applied to the data before the run, and it can be checked by running the
suite.

The fourth is the systems framing (Chapter~\ref{ch:systems} and
Section~\ref{sec:disc-systems}). It treats capacity as a stock observed once, and
the score, once used, as part of a loop with the service that acts on it. From
this it derives requirements that no measurement property can meet on its own:
report the domain scores beside the combined score, state the decision rule with
the score, monitor calibration over time and by equity stratum, and record the
actions taken. These are hypotheses and design requirements, not findings, and
they can be checked against implementation studies.

The fifth is the translation framework for haemophilia
(Chapter~\ref{ch:translational}). It keeps the three domains and replaces their
indicators with haemophilia-specific ones, through expert consensus,
psychometric re-estimation, recalibration and external validation, in that order.
It is a method for building a disease-specific version, not a validated
haemophilia measure, and it ends at a measure, not at a treatment recommendation.

The pre-registered reading rules of Section~\ref{sec:disc-reading} belong with
these. The thesis fixes the reading of each possible result as well as the
analysis, and that is what lets the contributions below be stated now without
waiting to see which way the results fall.

\paragraph{Contributions that depend on the NHATS results.}
If the measurement results fall on the supported branch, they would establish
that health capacity with three distinct domains can be measured from the
indicators of a general ageing survey, with adequate fit, reliability and, where
tested, invariance across sex and age band (Chapter~\ref{ch:results1}).

If the prognostic results fall on the supported branch, they would establish that
this score, built without reference to any outcome, predicts five-year mortality
and, separately, a first overnight hospital stay, with discrimination and
calibration that survive internal validation (Chapter~\ref{ch:results2}).

If the comparison falls on the supported branch, it would establish that the
score adds information, at decision thresholds in use, that the frailty index,
the QALY proxy and comorbidity measures do not supply, and that the gain survives
removal of the indicators the score shares with the frailty index. This is the
result that would justify a service collecting the extra information.

If the head-to-head test falls on the supported branch, it would establish that
the non-compensatory realisation model predicts better than an additive composite
of the same domains, which would give the weakest-domain reading of
Section~\ref{sec:disc-sys-weakest} empirical support in this cohort.

Each of these four would hold for one development cohort and would still need
external validation. If a result falls on another branch, the corresponding claim
is narrowed as Table~\ref{tab:disc-patterns} and
Figure~\ref{fig:disc-decision} fix in advance, and the five contributions that
stand now do not depend on it.

% =====================================================================
\section{What This Chapter Settles}
\label{sec:disc-close}
% =====================================================================

\noindent The chapter began by asking what the completed analyses would entitle
anyone to say. Part of the answer can be given now, and part must wait.

What it settles does not depend on the numbers. It fixes, before they arrive, how
each pattern of results will be read (Section~\ref{sec:disc-reading} and
Table~\ref{tab:disc-patterns}), and it names the results that would count against
the framework, so that a weak result cannot be explained away after the fact. It
places the construct against frailty indices, utility measures and the social
determinants literature, and states what a gain over each would and would not show,
including that the comparators are versions built from NHATS items. It sets each
strength beside the limitation that comes from the same decision, and carries into
the discussion the shortcomings that the checks of the data against their sources
left open: approximate event times, an extrapolation for proxy-reported cognition,
standard errors for income that are somewhat too small, and a sample that excludes
people in residential care. And it reads the thesis as a systems problem, which
shows that capacity under the realisation function is a property of the whole, that
a change does most in the weakest domain to a degree set by a parameter the data
cannot fix, and that the most serious risks of using the score lie in the rule
attached to it and the feedback its users receive, not in the score.

What it leaves open is everything that waits on the placeholders: whether the three
domains hold, whether the score predicts either outcome, whether it adds to the
frailty index once overlap is removed, and whether the conclusions survive the
floor of the domain scale. It also leaves open every systems claim, since none of
the loops, delays or control effects has been observed.

\noindent This chapter has considered the broader significance of the research by examining its implications, identifying its principal strengths and limitations, and outlining directions for future investigation. The thesis has argued that health may be understood as a multidimensional latent capability emerging through interactions among biological reserve, social resources and healthcare systems. Building upon this conceptual foundation, a structured methodology was developed for measuring the construct, evaluating its prognostic performance and translating it to disease-specific settings.

The implications of this work extend beyond the proposed $H_{\text{potential}}$ itself. Conceptually, the research contributes to ongoing discussions regarding the measurement of healthy ageing, resilience and capability. Methodologically, it demonstrates a staged framework linking theory development, psychometric evaluation and prognostic modelling while reducing analytical circularity. Translationally, it establishes a systematic pathway through which disease-independent health constructs may be adapted to specialised clinical populations without sacrificing conceptual coherence. At the same time, the conclusions of the thesis remain appropriately bounded: development and internal validation within a single cohort provide an essential foundation but do not constitute definitive evidence of universal generalisability, and external validation, recalibration and prospective clinical evaluation therefore represent necessary next steps before routine application can be considered.

Taken together, the discussion demonstrates that the $H_{\text{potential}}$ should be viewed as a research framework that has progressed from theoretical concept to empirically testable construct. Whether future studies ultimately confirm, refine or extend the present findings, the framework provides a coherent foundation for investigating health as a dynamic, multidimensional capability rather than solely a description of current disease. The next chapter draws together the principal conclusions of the thesis and explicitly states the original contributions made by this programme of research.

That chapter is needed because a discussion of this kind spreads its claims across
many sections and qualifies each one where it arises. A reader can finish it knowing
what every result would mean and still not have in one place what the thesis claims
as its own. Chapter~\ref{ch:conclusions} supplies that: the contribution, an answer
to each aim, and the boundary of each claim, stated once.

% =====================================================================
% LEFTOVER DRAFT, September 2026.
% Everything between \iffalse and \fi below duplicates the sections above
% (Summary of Principal Findings, Interpretation, Strengths, Limitations,
% Implications, Future Work, Summary) in an earlier and shorter form.
% It is fenced off rather than deleted so nothing is lost. The longer
% versions above supersede it. Delete this block when you are satisfied.
% =====================================================================
\iffalse
% =====================================================================
% NOTE TO AUTHOR: everything from here to the end of the file duplicates the
% sections above (Summary of Principal Findings, Interpretation, Strengths,
% Limitations, Implications, Future Work, Summary) in an earlier, shorter draft.
% It is left untouched apart from spelling. Delete it, or delete the longer
% versions above, before compiling; as it stands the chapter has two of each
% section and the duplicated headings will confuse the table of contents.
% =====================================================================
\section{Summary of Principal Findings}
\noindent This study set out to develop, validate, and test a multidimensional health metric, the $H_{\text{potential}}$, and to assess its translational relevance to haemophilia. The principal findings correspond to the four aims: the conceptual framework of Chapter~\ref{ch:theory}; the measurement model of Chapter~\ref{ch:results1}; the prognostic evaluation of Chapter~\ref{ch:results2}; and the translational mapping of Chapter~\ref{ch:translational}.

\medskip
\noindent\emph{[Placeholder: a concise statement of what was found. Whether the $H_{\text{potential}}$ formed a valid higher-order construct; whether it predicted five-year mortality and hospitalisation; whether it added value over conventional measures; and the standing of the haemophilia mapping.]}

\section{Interpretation in the Context of the Literature}
\noindent The $H_{\text{potential}}$ was designed to close the biomedical--social gap left by existing measures. QALYs and DALYs summarise burden for economic evaluation, quality-of-life instruments describe current status \citep{whitehead_2010}, and frailty indices quantify physiological reserve \citep{searle_2008}; none integrates biological, social, and structural determinants into a single predictive construct grounded in a theory of health \citep{sen_1993, who_2015_ageing, marmot_2005}. This section interprets how the results sit against that body of work: whether the three-domain structure held empirically, how the construct related to the comparators, and what the residual differences imply.

\medskip
\noindent\emph{[Placeholder: interpretation of the measurement and prognostic results against the literature, including how the $H_{\text{potential}}$ compares with frailty, QALY-based, and comorbidity measures, and what any incremental value means.]}

\section{Strengths}
\noindent The study has several strengths that hold regardless of the direction of the results.
\begin{itemize}
  \item \textbf{Theory-grounded construct.} The three domains are anchored in the capability approach, healthy ageing, frailty science, and the social determinants of health, rather than assembled for statistical convenience.
  \item \textbf{Separation of measurement and prediction.} Deriving the construct by structural equation modelling before, and independently of, the prognostic analysis reduces the circularity that affects composite indices whose weights are fitted to the outcomes they later predict.
  \item \textbf{Integration of social and structural determinants.} The framework treats social and systemic factors as intrinsic components of the health construct, not as adjustment covariates added after the fact.
  \item \textbf{Standardised, transparent reporting.} The study follows recognised standards for prediction models and measurement properties \citep{collins_2024, moons_2019}, with a pre-specified analysis plan and internal validation.
  \item \textbf{A defined translational path.} The haemophilia mapping and external-validation roadmap of Chapter~\ref{ch:translational} give a concrete route from a disease-independent metric to disease-specific use.
\end{itemize}

\section{Limitations}
\noindent The study also has limitations that qualify its conclusions.
\begin{itemize}
  \item \textbf{Development in a single cohort.} The metric is developed and internally validated in the NHATS cohort of community-dwelling United States adults aged 65 and older. Generalisability to other settings, age groups, and health systems is not established and requires external validation.
  \item \textbf{Internal, not external, validation.} Performance is assessed by bootstrap internal validation. External validation in an independent cohort is a distinct and necessary later phase.
  \item \textbf{Observational prognostic analysis.} The prognostic models are associational. Residual confounding cannot be excluded, and the hazard ratios should not be read as causal effects.
  \item \textbf{Measurement error carried forward.} The $H_{\text{potential}}$ enters the prognostic model as an estimated factor score. Although the bootstrap refits the full pipeline to reflect this, the two-stage design still carries measurement uncertainty into the prediction stage.
  \item \textbf{Indicator availability.} The domains are operationalised with the indicators NHATS provides. Some theoretically relevant constructs may be measured only in part, which can affect the content validity of a domain.
  \item \textbf{Assumptions in missing-data handling.} Multiple imputation assumes data are missing at random. Departures from this assumption would bias the estimates.
  \item \textbf{Normative choices.} Reverse-coding, standardisation, and the treatment of the systemic burden domain embed normative choices about how disadvantage should count. These are made explicit but remain open to reasoned disagreement.
\end{itemize}

\section{Implications}
\noindent Subject to the results, the study has implications at three levels.
\begin{itemize}
  \item \textbf{For measurement.} It offers a way to represent health as future capability rather than current status, integrating determinants that existing metrics treat separately.
  \item \textbf{For research.} The two-stage design provides a template for building theory-grounded composite measures and testing their prognostic value without circularity, extendable to other chronic conditions.
  \item \textbf{For practice and policy.} Any application to risk stratification or resource decisions should follow external validation and recalibration in the target population. The metric is a research instrument at this stage, not a validated clinical or policy tool.
\end{itemize}

\medskip
\noindent\emph{[Placeholder: sharpen the implications once the results are known, keeping claims within what internal validation in a single cohort supports.]}

\section{Future Work}
\noindent Several directions follow from this work.
\begin{itemize}
  \item \textbf{External validation.} The immediate next phase is external validation of the $H_{\text{potential}}$ in independent cohorts, and recalibration where needed.
  \item \textbf{Disease-specific adaptation.} The haemophilia roadmap of Chapter~\ref{ch:translational} is the first disease-specific application, beginning with a South African haemophilia cohort, and the approach extends to other chronic conditions of ageing and multimorbidity.
  \item \textbf{Dynamical extension.} The static construct can be generalised to a time-varying form. The dynamical ODE--SDE--PDE architecture set out in Appendix~\ref{app:dynamical}, and its haemophilia specialisation in Appendix~\ref{app:hemo_dynamics}, are directions for future modelling, to be developed and validated separately from the measurement-and-prediction work reported here.
  \item \textbf{Longer-term methodological extensions.} Later work could explore inverse modelling and Bayesian calibration to estimate the construct directly from large longitudinal data, and control-theoretic approaches to time-dependent intervention. These are theoretical directions and are not claims of the present thesis.
\end{itemize}

\section{Summary}
\noindent This chapter interpreted the thesis as a whole and set out its strengths, limitations, and implications. The next chapter draws the conclusions and states the contribution of the work.

\fi
